%\bibliographystyle{apsrev}
%\input{tcilatex}


\documentclass[aps,preprint,showpacs,showkeys]{revtex4}
%%%%%%%%%%%%%%%%%%%%%%%%%%%%%%%%%%%%%%%%%%%%%%%%%%%%%%%%%%%%%%%%%%%%%%%%%%%%%%%%%%%%%%%%%%%%%%%%%%%%%%%%%%%%%%%%%%%%%%%%%%%%
\usepackage{amsfonts}
\usepackage{amsmath}
\usepackage[tcidvi]{graphicx}

\setcounter{MaxMatrixCols}{10}
%TCIDATA{OutputFilter=LATEX.DLL}
%TCIDATA{Version=4.00.0.2321}
%TCIDATA{Created=Saturday, December 29, 2001 12:06:49}
%TCIDATA{LastRevised=Wednesday, July 31, 2002 13:18:31}
%TCIDATA{<META NAME="GraphicsSave" CONTENT="32">}
%TCIDATA{<META NAME="DocumentShell" CONTENT="Articles\SW\REVTeX 4">}
%TCIDATA{Language=American English}
%TCIDATA{CSTFile=revtex4.cst}

\newtheorem{theorem}{Theorem}
\newtheorem{acknowledgement}[theorem]{Acknowledgement}
\newtheorem{algorithm}[theorem]{Algorithm}
\newtheorem{axiom}[theorem]{Axiom}
\newtheorem{claim}[theorem]{Claim}
\newtheorem{conclusion}[theorem]{Conclusion}
\newtheorem{condition}[theorem]{Condition}
\newtheorem{conjecture}[theorem]{Conjecture}
\newtheorem{corollary}[theorem]{Corollary}
\newtheorem{criterion}[theorem]{Criterion}
\newtheorem{definition}[theorem]{Definition}
\newtheorem{example}[theorem]{Example}
\newtheorem{exercise}[theorem]{Exercise}
\newtheorem{lemma}[theorem]{Lemma}
\newtheorem{notation}[theorem]{Notation}
\newtheorem{problem}[theorem]{Problem}
\newtheorem{proposition}[theorem]{Proposition}
\newtheorem{remark}[theorem]{Remark}
\newtheorem{solution}[theorem]{Solution}
\newtheorem{summary}[theorem]{Summary}
\newenvironment{proof}[1][Proof]{\noindent\textbf{#1.} }{\ \rule{0.5em}{0.5em}}
\input{tcilatex}

\begin{document}

\preprint{AGPOPT}
\title[Optimized AGP]{An Effective Procedure for Energy Optimizing
Antisymmetrized Geminal Power States}
\author{B. Weiner}
\affiliation{Department of Physics, Pennsylvania State University, DuBois PA 15801}
\author{J. V. Ortiz}
\affiliation{Department of Chemistry, Kansas State University, Manhattan, KS 66506-3701}
\keywords{AGP, Correlation, Self Consistent Calculations}

\begin{abstract}
A procedure for energy optimizing an Antisymmetrized Geminal Power State,
which contains a description of correlation effects in an unbiased fashion,
is presented.

This procedure overcomes difficulties and shortcomings of past optimization
procedures for Antisymmetrized Geminal Power states. It is shown that the
computational cost scales as $r^{5}$, where $r$ is the size of the spin
orbital atomic basis set, which is superior to the scaling costs of multi
excited Complete Active Space Self Consistent Field calculations, but unlike
such calculations their is no preselection of configurations or excitation
level.

The variational parameters are the geminal coefficients with respect to a 
\emph{fixed }atomic orbital basis, which allows one to classify these
variables into chemically significant and non significant ones according to
their magnitudes and thus reduce size of calculations.
\end{abstract}

\date{05/02/2002}
\accepted{06/20/2002}
\received{05/09/2002}
\published{09/15/2002}
\startpage{1}
\maketitle

%\pacs{31.15.A, 03.65.F, 31.15.N, 31.25}

\section{Introduction\label{Introduction}}

Independent Particle States (IPS) have proved a mainstay of molecular
multi-particle quantum mechanics both in physical modelling and numerical
applications for many decades. The major attraction of these states lies in
the conceptually concrete and chemically understandable pictures of
molecular structure that they provide. Apart from an initial choice of $s$
atomic orbitals and possibly a decision about the molecular point group and
spin symmetry, one can obtain an unbiased description of the electronic
states of a molecule by finding an IPS\ that minimizes the energy. These
calculations scale at worst as $r^{4}$ where $r=2s$, and it has been shown
that they asymptotically scale as $r^{2.6}$. A major drawback of IPS\ models
is that they do not describe electron - electron correlation effects, which
can be crucial in correctly predicting the chemical and physical properties
of the system under study. Methods that introduce these effects, in an ab
initio manner, invariably increase the computational cost of any
calculation, scale at best as $r^{5}$, and usually involve constraints on
the molecular state description that reflect some bias beyond the choice of
the atomic orbital basis.

There is, however, another state ansatz that preserves most of the benefits
of IPS's, in particular providing an unbiased description of electronic
states, while at the same time introducing a description of correlation
effects. These are the Antisymmetrized Geminal Powers (AGP) states that were
first introduced to quantum chemistry by Coleman \cite%
{Coleman63,Coleman65,Coleman72} in the context of the analysis of the
N-representability problem. The class of AGP\ states in fact includes the
class of IPS's as a special case and so presents a true generalization of
the independent particle ansatz. In the following decade investigators \cite%
{Kutzenlig64, Silver69,Silver70,Silver71} examined the use of such states
and Bratoz and Durand \cite{Bratoz65,Bratoz69} constructed an energy
optimization procedure. However this procedure was not numerically
competitive with Hartree-Fock Self Consistent Field Methods nor with methods
that described correlation by other means. In the late 1970's it was found 
\cite{Lindeberg77,Lindeberg79,Lindeberg80} that AGP states produce the best
possible self consistent reference state for the Random Phase Approximation
(RPA)\ with exchange. This lead to a series of calculations, using real
restricted AGPs, on diatomic molecules that produced potential energy curves
for several states simultaneously and the oscillator strengths for
transitions amongst them over a range of internuclear distances from inside
the equilibrium position to the dissociation limit \cite%
{Kurtz83,Weiner84a,Weiner84,Weiner85,Weiner85b,Weiner87,Sangfelt87,Edwards84,Jensen82,Jorgen83}%
. The results of these calculations were in overall good agreement with
experiment, however the method was plagued with numerical problems that
became worse when it was applied to larger molecules.

In this paper we present a parametrization of AGP\ states and an
optimization procedure that overcomes these previously encountered numerical
difficulties, which at the same time has computational costs that increases
more slowly, with respect to the scaling with size of the basis set, than
comparable CASSCF and MCSCF procedures. The form of the parametrization
presented allows one to carry out all the calculations in a \emph{fixed}
localized atomic orbital basis and, as shown in the Appendix \ref{CostofComp}%
, scales as a $r^{5}$, where $r$ is the size of the spin orbital basis set.
This scaling should be compared to the \emph{combined }cost of the orbital
variation (i.e. integral transformations) and C.I variation steps in a
CASSCF, when almost all orders of excitation are described, which scale in a
worse fashion due to the large number of independent CI coefficients
contained in such states.

The AGP\ state is determined by a geminal, which in a given one particle
basis is determined by $r\left( r-1\right) /2$ complex variables, in
comparison an IPS\ is determined by $Nr$ complex variables where $N$ is the
number of electrons. These complex variables act as coordinates of the
nonlinear manifold of AGP\ states and IPS's respectively and the energy
associated with states belonging to these manifolds can be expressed as a
function of these variables. In order to obtain the best approximation to
the ground state of a given physical system one needs to find a global
minimum of this energy function. Different choices of one particle bases
induce different coordinate systems and thus various forms for the energy
function and its derivatives. By choosing between these forms we construct a
numerical scheme for the energy optimization of AGP\ states that scales in
the manner described above. In order to consider the transformations between
these coordinate systems in full generality\ we first carefully review in
section \ref{Bases and Geminals} transformations between general spin
orbitals and atomic orbitals which determine transformations of the
coordinates labelling geminals. In section \ref{The AGP State} we then
examine the definition of the AGP state, various forms of the energy
expression for such\ states and the transformations between these forms. In
section \ref{First Derivatives} we carry out the analogous analysis for the
first derivatives and in section \ref{Methods of Optimization} construct an
optimization scheme using the results of sections \ref{The AGP State} and %
\ref{First Derivatives}.

\section{Bases and Geminals\label{Bases and Geminals}}

Multi-electronic states are invariably constructed in terms of one particle
functions $\left\{ \varphi _{j}:1\leq j\leq r\right\} $ of space and spin
coordinates $\left( \mathbf{r},\xi \right) $ that are products of real
localized atomic orbitals $\left\{ \chi _{j}:1\leq j\leq s\right\} $ and
spin functions $\left\{ \alpha ,\beta \right\} $ i.e.%
\begin{eqnarray}
\varphi _{j}\left( \mathbf{r},\xi \right) &=&\chi _{j}\left( \mathbf{r}%
\right) \alpha \left( \xi \right) ;\quad 1\leq j\leq s  \notag \\
\varphi _{j+s}\left( \mathbf{r},\xi \right) &=&\chi _{j}\left( \mathbf{r}%
\right) \beta \left( \xi \right) ;\quad s\leq j\leq r
\end{eqnarray}%
where $r=2s$. The functions $\left\{ \chi _{j}\right\} $ are linearly
independent and thus produce a non singular overlap matrix $\mathbb{S}$
given by 
\begin{equation}
S_{ij}=\left\langle \chi _{i}|\chi _{j}\right\rangle =\int \chi _{i}\left( 
\mathbf{r}\right) \chi _{j}\left( \mathbf{r}\right) d\mathbf{r;\quad }1\leq
i,j\leq s
\end{equation}%
The spin functions $\left\{ \alpha ,\beta \right\} $ are orthonormal i.e. 
\begin{eqnarray}
\left\langle \alpha |\alpha \right\rangle &=&\left\langle \beta |\beta
\right\rangle =1  \notag \\
\left\langle \alpha |\beta \right\rangle &=&\left\langle \beta |\alpha
\right\rangle =0
\end{eqnarray}%
where spin inner product is 
\begin{equation}
\left\langle \sigma |\sigma ^{\prime }\right\rangle =\int \sigma ^{\ast
}\left( \xi \right) \sigma ^{\prime }\left( \xi \right) d\xi
\end{equation}%
and $^{\ast }$ denotes complex conjugation. The functions $\left\{ \varphi
_{j};1\leq j\leq r\right\} $ form the one particle Hilbert space $\mathcal{H}%
^{1}$. Orthonormal molecular orbitals $\left\{ \omega _{j}:1\leq j\leq
s\right\} $ can be defined in terms of the atomic orbitals as 
\begin{equation}
\left( \underline{\omega }\right) =\left( \underline{\chi }\right) \mathbb{W}
\end{equation}%
where we have used the array of basis functions notation 
\begin{eqnarray}
\left( \underline{\omega }\right) &=&\left( \omega _{1},\ldots ,\omega
_{s}\right)  \notag \\
\left( \underline{\chi }\right) &=&\left( \chi _{1},\ldots ,\chi _{s}\right)
\end{eqnarray}%
The linear transformation $\mathbb{W}$ has the properties 
\begin{equation}
\mathbb{W}^{\dagger }\mathbb{SW}=\mathbb{I}_{s}\,\mathbb{\Longleftrightarrow
WW}^{\dagger }=\mathbb{S}^{-1}\mathbb{\Longleftrightarrow W}^{-1}=\mathbb{W}%
^{\dagger }\mathbb{S}
\end{equation}%
and is clearly non unique as any transformation $\mathbb{WU}$ has the same
properties when $\mathbb{U}$ is unitary. The molecular orbitals $\left\{
\omega _{j}:1\leq j\leq s\right\} $ can be used to form \emph{restricted
spin }bases of the one particle space $\mathcal{H}^{1}$ given by 
\begin{eqnarray}
\psi _{j}\left( \mathbf{r},\xi \right) &=&\omega _{j}\left( \mathbf{r}%
\right) \alpha \left( \xi \right) ;\quad 1\leq j\leq s  \notag \\
\psi _{j+s}\left( \mathbf{r},\xi \right) &=&\omega _{j}\left( \mathbf{r}%
\right) \beta \left( \xi \right) ;\quad s<j\leq r
\end{eqnarray}%
so that 
\begin{equation}
\left( \underline{\psi }\right) =\left( \underline{\varphi }\right) \left( 
\begin{array}{cc}
\mathbb{W} & 0 \\ 
0 & \mathbb{W}%
\end{array}%
\right)
\end{equation}%
This construction can be generalized to form \emph{unrestricted spin }bases
for $\mathcal{H}^{1}$ given by 
\begin{equation}
\left( \underline{\psi }\right) =\left( \underline{\varphi }\right) \left( 
\begin{array}{cc}
\mathbb{W}_{\alpha } & 0 \\ 
0 & \mathbb{W}_{\beta }%
\end{array}%
\right)
\end{equation}%
where $\mathbb{W}_{\alpha }\neq \mathbb{W}_{\beta }$ but 
\begin{equation}
\mathbb{W}_{\alpha }^{\dagger }\mathbb{SW}_{\alpha }=\mathbb{W}_{\beta
}^{\dagger }\mathbb{SW}_{\beta }=\mathbb{I}_{s}\,
\end{equation}%
One can extend this type of construction further to form \emph{generalized
spin bases }%
\begin{equation}
\left( \underline{\psi }\right) =\left( \underline{\varphi }\right) \left( 
\begin{array}{cc}
\mathbb{W}_{\alpha \alpha } & \mathbb{W}_{\alpha \beta } \\ 
\mathbb{W}_{\beta \alpha } & \mathbb{W}_{\beta \beta }%
\end{array}%
\right)
\end{equation}%
where $\mathbb{W}_{\alpha \beta }\neq 0$ and $\mathbb{W}_{\beta \alpha }\neq
0$ but 
\begin{equation}
\left( 
\begin{array}{cc}
\mathbb{W}_{\alpha \alpha } & \mathbb{W}_{\alpha \beta } \\ 
\mathbb{W}_{\beta \alpha } & \mathbb{W}_{\beta \beta }%
\end{array}%
\right) ^{\dagger }\left( 
\begin{array}{cc}
\mathbb{S} & 0 \\ 
0 & \mathbb{S}%
\end{array}%
\right) \left( 
\begin{array}{cc}
\mathbb{W}_{\alpha \alpha } & \mathbb{W}_{\alpha \beta } \\ 
\mathbb{W}_{\beta \alpha } & \mathbb{W}_{\beta \beta }%
\end{array}%
\right) =\left( 
\begin{array}{cc}
\mathbb{I}_{s} & 0 \\ 
0 & \mathbb{I}_{s}%
\end{array}%
\right) ,
\end{equation}%
\begin{equation}
\left( 
\begin{array}{cc}
\mathbb{W}_{\alpha \alpha } & \mathbb{W}_{\alpha \beta } \\ 
\mathbb{W}_{\beta \alpha } & \mathbb{W}_{\beta \beta }%
\end{array}%
\right) \left( 
\begin{array}{cc}
\mathbb{W}_{\alpha \alpha } & \mathbb{W}_{\alpha \beta } \\ 
\mathbb{W}_{\beta \alpha } & \mathbb{W}_{\beta \beta }%
\end{array}%
\right) ^{\dagger }=\left( 
\begin{array}{cc}
\mathbb{S} & 0 \\ 
0 & \mathbb{S}%
\end{array}%
\right) ^{-1}
\end{equation}%
and 
\begin{equation}
\left( 
\begin{array}{cc}
\mathbb{W}_{\alpha \alpha } & \mathbb{W}_{\alpha \beta } \\ 
\mathbb{W}_{\beta \alpha } & \mathbb{W}_{\beta \beta }%
\end{array}%
\right) ^{-1}=\left( 
\begin{array}{cc}
\mathbb{W}_{\alpha \alpha } & \mathbb{W}_{\alpha \beta } \\ 
\mathbb{W}_{\beta \alpha } & \mathbb{W}_{\beta \beta }%
\end{array}%
\right) ^{\dagger }\left( 
\begin{array}{cc}
\mathbb{S} & 0 \\ 
0 & \mathbb{S}%
\end{array}%
\right)
\end{equation}%
which we can write in non partitioned form as 
\begin{equation}
\mathbb{W}_{f}^{-1}=\mathbb{W}_{f}^{\dagger }\mathbb{S}_{f}
\end{equation}%
where the subscript ''$f$'' denotes full and 
\begin{equation}
\mathbb{A}_{f}\equiv \left( 
\begin{array}{cc}
\mathbb{A}_{\alpha \alpha } & \mathbb{A}_{\alpha \beta } \\ 
\mathbb{A}_{\beta \alpha } & \mathbb{A}_{\beta \beta }%
\end{array}%
\right) .
\end{equation}%
The individual basis functions in the \emph{generalized spin bases }are of
the form 
\begin{eqnarray}
\left| \psi _{j}\right\rangle &=&\sum_{1\leq k\leq s}\left\{ \left| \varphi
_{k}\right\rangle W_{\alpha \alpha \,kj}+\left| \varphi _{k+s}\right\rangle
W_{\beta \alpha \,kj}\right\} \equiv \left| \psi _{j}^{\alpha }\right\rangle
+\left| \psi _{j}^{\beta }\right\rangle ;\quad 1\leq j\leq s  \notag \\
\left| \psi _{j+s}\right\rangle &=&\sum_{1\leq k\leq s}\left\{ \left|
\varphi _{k}\right\rangle W_{\alpha \beta \,kj}+\left| \varphi
_{k+s}\right\rangle W_{\beta \beta \,kj}\right\} \equiv \left| \psi
_{j+s}^{\alpha }\right\rangle +\left| \psi _{j+s}^{\beta }\right\rangle
;\quad 1\leq j\leq s
\end{eqnarray}

Geminals, which are antisymmetric two particle functions, can be expanded in
terms of the two particle orthonormal basis $\left\{ \left| \psi _{j}\psi
_{k}\right\rangle ;1\leq j<k\leq r\right\} $ formed from antisymmetrized
products of the generalized spin orbitals $\left\{ \left| \psi
_{j}\right\rangle ;1\leq j\leq r\right\} $ as 
\begin{equation}
\left| g\right\rangle =\sum_{1\leq j<k\leq r}c_{jk}\left| \psi _{j}\psi
_{k}\right\rangle
\end{equation}%
The same geminal $\left| g\right\rangle $ can also be expanded in terms of
the non orthonormal two particle basis, $\left\{ \left| \varphi _{j}\varphi
_{k}\right\rangle ;1\leq j<k\leq r\right\} $, as 
\begin{equation}
\left| g\right\rangle =\sum_{1\leq j<k\leq r}\gamma _{jk}\left| \varphi
_{j}\varphi _{k}\right\rangle
\end{equation}%
The sets of coefficients $\left\{ c_{jk}\right\} $ and $\left\{ \gamma
_{jk}\right\} $ are just coordinates for the same geminal in different
coordinate systems. The transformations between the coordinate systems $%
\mathbf{c}$ and $\mathbf{\gamma }$ are \emph{linear }and given by (see
Appendix \ref{appengct})%
\begin{eqnarray}
\mathbf{\gamma } &=&\,\left( \mathbb{W}_{f}\otimes \mathbb{W}_{f}\right) 
\mathbf{c=\,}\mathbb{W}_{f}\mathbf{c}\mathbb{W}_{f}^{t}  \notag \\
\mathbf{c} &=&\left( \mathbb{W}_{f}\otimes \mathbb{W}_{f}\right) ^{-1}\,%
\mathbf{\gamma =\,}\mathbb{W}_{f}^{-1}\,\mathbf{\gamma \,}\mathbb{W}%
_{f}^{-1t}=\mathbf{\,}\mathbb{W}_{f}^{\dagger }\mathbb{S}_{f}\,\,\mathbf{%
\gamma \,\,}\mathbb{S}_{f}^{t}\mathbb{W}_{f}^{\ast }
\end{eqnarray}%
where $\mathbf{c}$ and $\mathbf{\gamma }$ are viewed first as vectors and
then matrices. As the coordinate transformation is linear one has the
following relationship between partial derivative w.r.t the variables $%
\mathbf{c}$ and $\mathbf{\gamma }$ 
\begin{eqnarray}
\frac{\partial }{\partial \mathbf{\gamma }} &=&\left( \mathbb{W}_{f}\otimes 
\mathbb{W}_{f}\right) \frac{\partial }{\partial \mathbf{c}}\equiv \mathbb{W}%
_{f}\frac{\partial }{\partial \mathbf{c}}\mathbb{W}_{f}^{t}  \notag \\
\frac{\partial }{\partial \mathbf{c}} &=&\left( \mathbb{W}_{f}\otimes 
\mathbb{W}_{f}\right) ^{-1}\,\frac{\partial }{\partial \mathbf{\gamma }}%
\equiv \mathbf{\,\,}\mathbb{W}_{f}^{\dagger }\mathbb{S}_{f}\,\,\frac{%
\partial }{\partial \mathbf{\gamma }}\mathbf{\,\,}\mathbb{S}_{f}^{t}\mathbb{W%
}_{f}^{\ast }  \label{coordtrans}
\end{eqnarray}%
where again $\frac{\partial }{\partial \mathbf{\gamma }}$ and $\frac{%
\partial }{\partial \mathbf{c}}$ are viewed first as vectors and then
matrices.

\section{The AGP State\label{The AGP State}}

\subsection{Definition}

We can associate a geminal creator, $G^{\dagger }$, with a geminal, $\left|
g\right\rangle $, by 
\begin{equation}
G^{\dagger }=\frac{1}{2}\sum_{1\leq j,k\leq r}\gamma _{jk}a_{j}^{\dagger
}a_{k}^{\dagger }=\frac{1}{2}\sum_{1\leq j,k\leq r}c_{jk}b_{j}^{\dagger
}b_{k}^{\dagger }
\end{equation}%
where $\gamma _{jk}=-\gamma _{kj}$ and $c_{jk}=-c_{kj}$ and the second
quantized creators have the properties 
\begin{eqnarray}
a_{j}^{\dagger }\left| \phi \right\rangle &=&\left| \varphi _{j}\right\rangle
\notag \\
b_{j}^{\dagger }\left| \phi \right\rangle &=&\left| \psi _{j}\right\rangle
\end{eqnarray}%
when acting on the zero particle state $\left| \phi \right\rangle $ and
satisfy the commutation relationships 
\begin{eqnarray}
\left[ a_{j},a_{k}^{\dagger }\right] _{+} &=&S_{jk}  \notag \\
\left[ b_{j},b_{k}^{\dagger }\right] _{+} &=&\delta _{jk}  \notag \\
\left[ a_{j},a_{k}\right] _{+} &=&\left[ a_{j}^{\dagger },a_{k}^{\dagger }%
\right] _{+}=\left[ b_{j},b_{k}\right] _{+}=\left[ b_{j}^{\dagger
},b_{k}^{\dagger }\right] _{+}=0
\end{eqnarray}%
The geminal can then be expressed in second quantized, \emph{coordinate free}%
, notation as 
\begin{equation}
\left| g\right\rangle =G^{\dagger }\left| \phi \right\rangle
\end{equation}%
\qquad In order to display the coordinate of the geminal explicitly we will
use the notation 
\begin{eqnarray}
\left| g\left( \mathbf{\gamma }\right) \right\rangle &=&G\left( \mathbf{%
\gamma }\right) ^{\dagger }\left| \phi \right\rangle =\frac{1}{2}\sum_{1\leq
j,k\leq r}\gamma _{jk}a_{j}^{\dagger }a_{k}^{\dagger }\left| \phi
\right\rangle  \notag \\
\left| g\left( \mathbf{c}\right) \right\rangle &=&G\left( \mathbf{c}\right)
^{\dagger }\left| \phi \right\rangle =\frac{1}{2}\sum_{1\leq j,k\leq
r}c_{jk}b_{j}^{\dagger }b_{k}^{\dagger }\left| \phi \right\rangle
\label{quantgem}
\end{eqnarray}%
which leads to the coordinate dependent expressions for the geminal 
\begin{equation}
\left| g\right\rangle =G\left( \mathbf{\gamma }\right) ^{\dagger }\left|
\phi \right\rangle =G\left( \mathbf{c}\right) ^{\dagger }\left| \phi
\right\rangle
\end{equation}

The 2N-electron \emph{unnormalized }AGP state, $\left\vert
g^{N}\right\rangle $, is defined as 
\begin{equation}
\left\vert g^{N}\right\rangle =\left( G^{\dagger }\right) ^{N}\left\vert
\phi \right\rangle
\end{equation}%
which can be expanded in coordinates wrt a linearly dependent \emph{unordered%
} bases as 
\begin{eqnarray}
\left\vert g^{N}\right\rangle &=&\frac{1}{2^{N}}\sum_{1\leq
j_{1},k_{1},\cdots ,j_{N},k_{N}\leq r}c_{j_{1}k_{1}}\cdots
c_{j_{N}k_{N}}\left\vert \psi _{j_{1}}\psi _{k_{1}}\cdots \psi _{j_{N}}\psi
_{k_{N}}\right\rangle  \notag \\
&=&\frac{1}{2^{N}}\sum_{1\leq j_{1},k_{1},\cdots ,j_{N},k_{N}\leq r}\gamma
_{j_{1}k_{1}}\cdots \gamma _{j_{N}k_{N}}\left\vert \varphi _{j_{1}}\varphi
_{k_{1}}\cdots \varphi _{j_{N}}\varphi _{k_{N}}\right\rangle .  \label{AGPEX}
\end{eqnarray}%
The sums in eq. $\left( \ref{AGPEX}\right) $ can be expressed wrt a linearly
independent \emph{ordered} basis by collecting together CI coefficients that
refer to configurations that can permuted into each other as described in
Appendix \ref{Pfaff}. The use of these ordered coefficients, however,
increases the complexity of the expressions and the construction of
optimization procedures.

\subsection{State Energy}

The energy of an 2N-electron AGP state is given by 
\begin{equation}
\mathcal{E}\left( g^{N}\right) =\frac{\mathcal{H}\left( g^{N}\right) }{%
\mathcal{N}\left( g^{N}\right) }  \label{Energy}
\end{equation}%
where%
\begin{equation}
\mathcal{H}\left( g^{N}\right) =\left\langle g^{N}|Hg^{N}\right\rangle
\end{equation}%
in terms of the Hamiltonian%
\begin{equation*}
H=h_{0}+h_{1}+h_{2}\qquad \qquad \qquad \qquad \qquad \qquad \qquad \qquad
\end{equation*}%
\begin{equation*}
h_{0}=\sum_{1\leq j<k\leq N_{nuc}}\frac{Z_{j}Z_{k}}{\left\Vert \mathbf{R}%
_{j}-\mathbf{R}_{k}\right\Vert }\qquad \qquad \qquad \qquad \qquad \quad
\end{equation*}%
\begin{equation*}
h_{1}=-\left( \sum_{_{1\leq j\leq N_{nuc};1\leq k\leq 2N}}\frac{Z_{j}}{%
\left\Vert \mathbf{R}_{j}-\mathbf{r}_{k}\right\Vert }+\frac{1}{2m_{e}}%
\sum_{1\leq j\leq 2N}\nabla _{\mathbf{r}_{j}}^{2}\right) \,
\end{equation*}%
\begin{equation}
h_{2}=\sum_{1\leq j<k\leq N}\frac{1}{\left\Vert \mathbf{r}_{j}-\mathbf{r}%
_{k}\right\Vert }\qquad \qquad \qquad \qquad \qquad \qquad \quad
\end{equation}%
and the normalization factor 
\begin{equation}
\mathcal{N}\left( g^{N}\right) =\left\langle g^{N}|g^{N}\right\rangle
\end{equation}%
for nuclei of charges $\left\{ Z_{j};1\leq j\leq N_{nuc}\right\} $ at a
fixed configuration $\left( \mathbf{R}_{1},\ldots ,\mathbf{R}%
_{N_{nuc}}\right) $. We can express the functions 
\begin{eqnarray}
\mathcal{E}\left( g^{N}\right) &=&\mathcal{E}\left( \mathbf{c}^{\ast },%
\mathbf{c}\right) =\mathcal{E}\left( \mathbf{\gamma }^{\ast },\mathbf{\gamma 
}\right)  \notag \\
\mathcal{H}\left( g^{N}\right) &=&\mathcal{H}\left( \mathbf{c}^{\ast },%
\mathbf{c}\right) =\mathcal{H}\left( \mathbf{\gamma }^{\ast },\mathbf{\gamma 
}\right)  \notag \\
\mathcal{N}\left( g^{N}\right) &=&\mathcal{N}\left( \mathbf{c}^{\ast },%
\mathbf{c}\right) =\mathcal{N}\left( \mathbf{\gamma }^{\ast },\mathbf{\gamma 
}\right)  \label{EnFunc}
\end{eqnarray}%
in terms of coordinates by observing%
\begin{eqnarray}
\mathcal{H}\left( \mathbf{c},\mathbf{c}^{\ast }\right) &=&\frac{1}{2^{2N}}%
\sum c_{j_{1}^{\prime }k_{1}^{\prime }}^{\ast }\cdots c_{j_{N}^{\prime
}k_{N}^{\prime }}^{\ast }c_{j_{1}k_{1}}\cdots c_{j_{N}k_{N}}\langle \psi
_{j_{1}^{\prime }}\psi _{k_{1}^{\prime }}\cdots \psi _{j_{N}^{\prime }}\psi
_{k_{N}^{\prime }}\left\vert H\psi _{j_{1}}\psi _{k_{1}}\cdots \psi
_{j_{N}}\psi _{k_{N}}\right\rangle  \notag \\
&&\qquad \qquad \qquad \qquad \quad _{1\leq j_{1},k_{1},j_{2},k_{2},\cdots
,j_{N},k_{N}\leq r;\,1\leq j_{1}^{\prime },k_{1}^{\prime },j_{2}^{\prime
},k_{2}^{\prime },\cdots ,j_{N}^{\prime },k_{N}^{\prime }\leq r}
\end{eqnarray}%
and 
\begin{eqnarray}
\mathcal{H}\left( \mathbf{\gamma }^{\ast },\mathbf{\gamma }\right) &=&\frac{1%
}{2^{2N}}\sum \gamma _{j_{1}^{\prime }k_{1}^{\prime }}^{\ast }\cdots \gamma
_{j_{N}^{\prime }k_{N}^{\prime }}^{\ast }\gamma _{j_{1}k_{1}}\cdots \gamma
_{j_{N}k_{N}}\left\langle \varphi _{j_{1}^{\prime }}\varphi _{k_{1}^{\prime
}}\cdots \varphi _{j_{N}^{\prime }}\varphi _{k_{N}^{\prime }}|H\varphi
_{j_{1}}\varphi _{k_{1}}\cdots \varphi _{j_{N}}\varphi _{k_{N}}\right\rangle
\notag \\
&&\qquad \qquad \qquad \qquad \quad _{1\leq j_{1},k_{1},j_{2},k_{2},\cdots
,j_{N},k_{N}\leq r;\,1\leq j_{1}^{\prime },k_{1}^{\prime },j_{2}^{\prime
},k_{2}^{\prime },\cdots ,j_{N}^{\prime },k_{N}^{\prime }\leq r}
\end{eqnarray}%
obtaining the norm functions $\mathcal{N}\left( \mathbf{c}^{\ast },\mathbf{c}%
\right) $ and $\mathcal{N}\left( \mathbf{\gamma }^{\ast },\mathbf{\gamma }%
\right) $ by setting $H=I$ in these expressions.

\subsection{Canonical Coordinates}

It is always possible to expand the AGP state in terms of $s$ complex
canonical geminal coordinates \cite{L�wdin56,Bloch62,Zumino62} $\left\{
\zeta _{j};1\leq j\leq s\right\} $, instead of the $s\left( 2s-1\right) $
complex coordinates used in eq. $\left( \ref{AGPEX}\right) $, as 
\begin{equation}
\left| g^{N}\right\rangle =N!\sum_{1\leq j_{1}<j_{2}<\cdots <j_{N}\leq
s}\zeta _{j_{1}}\cdots \zeta _{j_{N}}\left| \eta _{j_{1}}\eta
_{j_{1}+s}\cdots \eta _{j_{N}\,}\eta _{j_{N}+s}\right\rangle  \label{canexp}
\end{equation}%
where $\left\{ \eta _{j};1\leq j\leq r\right\} $ are canonical general spin
orbitals of the geminal $\left| g\right\rangle ,$ which has the canonical
expansion 
\begin{equation}
\left| g\right\rangle =\sum_{1\leq j\leq s}\zeta _{j}\left| \eta _{j}\eta
_{j+s}\right\rangle =\sum_{1\leq j\leq s}n_{j}^{\frac{1}{2}}e^{i\theta
_{j}}\left| \eta _{j}\eta _{j+s}\right\rangle ,
\end{equation}%
and we have expressed the canonical coordinates $\left\{ \zeta _{j}\right\} $
in polar form as 
\begin{equation}
\zeta _{j}=n_{j}^{\frac{1}{2}}e^{i\theta _{j}};\quad n_{j}=\left| \zeta
_{j}\right| ^{2}.
\end{equation}%
The canonical coordinates $\left\{ \zeta _{j}\right\} $ can be constructed
by first forming the one particle reduced density matrix $\frac{1}{2}\mathbf{%
cc}^{\dagger }$ of the two particle state $\left| g\right\rangle $
represented in the orthonormal basis $\left\{ \psi _{j};1\leq j\leq
r\right\} $ and then carrying out the following steps \cite{Bloch62,Zumino62}%
, which are explained in detail in Appendix $\ref{Appenccc}$:

\begin{enumerate}
\item diagonalize the matrix $\mathbf{cc}^{\dagger }$ formed from the
complex antisymmetric coefficient matrix $\mathbf{c}$ with a unitary
transformation $\mathbb{U}$ to obtain 
\begin{equation}
\mathbb{U}^{\dagger }\mathbf{cc}^{\dagger }\mathbb{U=}\left( 
\begin{array}{cc}
\mathbf{n} & \mathbf{0} \\ 
\mathbf{0} & \mathbf{n}%
\end{array}%
\right)  \label{uccu}
\end{equation}%
where 
\begin{equation}
\mathbf{n=}\left( 
\begin{array}{ccc}
n_{1} & 0 & 0 \\ 
0 & \ddots & 0 \\ 
0 & 0 & n_{s}%
\end{array}%
\right) ;\qquad n_{j}\geq 0,
\end{equation}%
by solving the eigenvalue problem to determine the columns $\left\{
u_{j}\right\} $ and eigenvalues $\left\{ n_{j}\right\} $ 
\begin{equation}
\mathbf{cc}^{\dagger }\mathbf{\,}u_{j}=n_{j}u_{j}
\end{equation}

\item then use $\mathbb{U}$ as a coordinate transformation matrix to obtain
a block diagonal matrix%
\begin{equation}
\mathbb{U}^{\dagger }\mathbf{c}\mathbb{U}^{\ast }=\left( 
\begin{array}{ccc}
\mathbb{C}_{1} & 0 & 0 \\ 
0 & \ddots & 0 \\ 
0 & 0 & \mathbb{C}_{p}%
\end{array}%
\right)
\end{equation}%
where the blocks $\left\{ \mathbb{C}_{\kappa },1\leq \kappa \leq p\right\} $
are antisymmetric, $d\left( \kappa \right) \times d\left( \kappa \right) $
in size, and have the property 
\begin{equation}
\mathbb{C}_{\kappa }\mathbb{C}_{\kappa }^{\dagger }=\eta _{\kappa }\mathbb{I}%
_{d\left( \kappa \right) };\quad 1\leq \kappa \leq p
\end{equation}%
and $\eta _{\kappa }$ is the value of the $d\left( \kappa \right) $-fold
degenerate $\kappa ^{th}$ eigenvalue of $\mathbf{cc}^{\dagger }$, then

\item simultaneously diagonalize the real and imaginary parts of $\mathbb{C}%
_{\kappa }$ with a real orthogonal $d\left( \kappa \right) \times d\left(
\kappa \right) $ matrix, which thus diagonalizes $\mathbb{C}_{\kappa }$,
(see Appendix $\ref{Appenccc}$), to produce 
\begin{equation}
\mathbb{O}_{\kappa }^{t}\mathbb{C}_{\kappa }\mathbb{O}_{\kappa }=\left( 
\begin{array}{cc}
\mathbf{0} & e^{i\mathbf{\phi }_{\kappa }} \\ 
-e^{i\mathbf{\phi }_{\kappa }} & \mathbf{0}%
\end{array}%
\right) ;\quad 1\leq \kappa \leq p  \label{orthog}
\end{equation}%
where 
\begin{equation}
e^{i\mathbf{\phi }_{\kappa }}=\left( 
\begin{array}{ccc}
e^{i\theta _{\kappa 1}} & 0 & 0 \\ 
0 & \ddots & 0 \\ 
0 & 0 & e^{i\theta _{\kappa d\left( \kappa \right) }}%
\end{array}%
\right)
\end{equation}%
hence producing the transformation 
\begin{equation}
\mathbb{V}^{t}\mathbf{c}\mathbb{V\equiv O}^{t}\mathbb{U}^{\dagger }\mathbf{c}%
\mathbb{U}^{\ast }\mathbb{O}=\left( 
\begin{array}{cc}
\mathbf{0} & \mathbf{\zeta } \\ 
-\mathbf{\zeta } & \mathbf{0}%
\end{array}%
\right)  \label{ucu}
\end{equation}%
where%
\begin{equation}
\mathbb{O}=\left( 
\begin{array}{ccc}
\mathbb{O}_{1} & 0 & 0 \\ 
0 & \ddots & 0 \\ 
0 & 0 & \mathbb{O}_{p}%
\end{array}%
\right)
\end{equation}%
and 
\begin{equation}
\mathbf{\zeta =}\left( 
\begin{array}{ccc}
\zeta _{1} & 0 & 0 \\ 
0 & \ddots & 0 \\ 
0 & 0 & \zeta _{s}%
\end{array}%
\right) .
\end{equation}
\end{enumerate}

The canonical coordinates can be constructed directly in the non orthogonal $%
\mathbf{\gamma }$ coordinate system by letting $\left\{ \psi _{j};1\leq
j\leq r\right\} $ be the orthonormal basis obtained from $\left\{ \varphi
_{j};1\leq j\leq r\right\} $ by symmetric orthonormalization which gives%
\begin{equation}
\mathbf{c}=\mathbb{S}_{f}^{\frac{1}{2}}\mathbf{\gamma }\mathbb{S}_{f}^{\frac{%
1}{2}t}
\end{equation}%
\begin{equation}
\mathbf{cc}^{\dagger }=\mathbb{S}_{f}^{\frac{1}{2}}\mathbf{\gamma }\mathbb{S}%
_{f}^{t}\mathbf{\gamma }^{\dagger }\mathbb{S}_{f}^{\frac{1}{2}}
\end{equation}%
and hence the following hermitian eigenvalue equation for the columns $%
\left\{ u_{j}\right\} $ and the eigenvalues $\left\{ n_{j}\right\} $ 
\begin{equation}
\mathbb{S}_{f}^{\frac{1}{2}}\mathbf{\gamma }\mathbb{S}_{f}^{t}\mathbf{\gamma 
}^{\dagger }\mathbb{S}_{f}^{\frac{1}{2}}\mathbf{\,}u_{j}=n_{j}\,u_{j};\qquad
1\leq j\leq r  \label{NGSO}
\end{equation}%
which produces the unitary transformation $\mathbb{U}$ that has the
properties 
\begin{equation}
\mathbb{U}^{\dagger }\,\mathbb{S}_{f}^{\frac{1}{2}}\mathbf{\gamma }\mathbb{S}%
_{f}^{t}\mathbf{\gamma }^{\dagger }\mathbb{S}_{f}^{\frac{1}{2}}\mathbf{\,}%
\mathbb{U}=\left( 
\begin{array}{cc}
\mathbf{n} & \mathbf{0} \\ 
\mathbf{0} & \mathbf{n}%
\end{array}%
\right) .
\end{equation}%
Then using $\mathbb{U}$ to block diagonalize $\mathbb{S}_{f}^{\frac{1}{2}}%
\mathbf{\gamma }\mathbb{S}_{f}^{\frac{1}{2}t}$, constructing the $\left\{ 
\mathbb{O}_{\kappa };1\leq \kappa \leq p\right\} $, as in eq. $\left( \ref%
{orthog}\right) $, then generating the transformation, as in eq. $\left( \ref%
{ucu}\right) $, that produces the skew canonical form, i.e.%
\begin{equation}
\mathbb{OU}^{\dagger }\mathbb{\,S}_{f}^{\frac{1}{2}}\mathbf{\gamma }\mathbb{S%
}_{f}^{\frac{1}{2}t}\mathbf{\,}\mathbb{U}^{\ast }\mathbb{O}^{t}\equiv 
\mathbb{V\,S}_{f}^{\frac{1}{2}}\mathbf{\gamma }\mathbb{S}_{f}^{\frac{1}{2}t}%
\mathbf{\,}\mathbb{V}^{t}=\left( 
\begin{array}{cc}
\mathbf{0} & \mathbf{\zeta } \\ 
-\mathbf{\zeta } & \mathbf{0}%
\end{array}%
\right) .  \label{eta}
\end{equation}

In general the canonical coordinate axes are very sensitive to changes in
the geminal. Small geminal perturbations therefore will make the right hand
side of eq. $\left( \ref{eta}\right) $ non skew diagonal if the axes, i.e.
basis functions, are not changed correspondingly.

\subsection{State Energy in Canonical Coordinates}

A. J. Coleman showed \cite{Coleman65} that the energy of a $2N$-electron
AGP\ state (see eq.$\left( \ref{Energy}\right) $ has a particularly succinct
form (see Appendix \ref{EnDetail}) when expressed in the canonical
coordinates $\left( \mathbf{\zeta }^{\ast },\mathbf{\zeta }\right) $. The
normalization $\mathcal{N}\left( g^{N}\right) $ becomes 
\begin{equation}
\mathcal{N}\left( \mathbf{\zeta }^{\ast },\mathbf{\zeta }\right) =\left(
N\right) !^{2}S_{N}
\end{equation}%
where $S_{N}$ is the $N^{th}$ order elementary symmetric function, which is
defined in terms of the eigenvalues $\left\{ n_{j};1\leq j\leq s\right\} $
of the one particle reduced density operator associated with the geminal $%
\left| g\right\rangle $ as 
\begin{equation}
S_{N}=\sum_{1\leq j_{1}<j_{2}<\cdots <j_{N}\leq s}n_{j_{1}}\cdots
n_{_{j_{N}}}.
\end{equation}%
The unnormalized energy $\mathcal{H}\left( g^{N}\right) $ in the canonical
coordinates $\left( \mathbf{\zeta }^{\ast },\mathbf{\zeta }\right) $ becomes%
\begin{eqnarray}
&&\mathcal{H}\left( \mathbf{\zeta }^{\ast },\mathbf{\zeta }\right) =h_{0}%
\mathcal{N}\left( g^{N}\right) +\left( N\right) !^{2}\left\{ \sum_{1\leq
j\leq s}n_{j}S_{N-1}\left( \hat{\jmath}\right) \left\{
h_{jj}+h_{j+sj+s}+\left\langle jj+s||jj+s\right\rangle \right\} \right. 
\notag \\
&&\qquad +\sum_{1\leq j\neq k\leq s}\zeta _{j}^{\ast }\zeta
_{k}S_{N-1}\left( \hat{\jmath}\hat{k}\right) \left\langle
jj+s||kk+s\right\rangle \left. +\sum_{1\leq j,k\leq
s}n_{j}n_{k}S_{N-2}\left( \hat{\jmath}\hat{k}\right) \left\langle
jk|||jk\right\rangle \right\}  \label{canen}
\end{eqnarray}%
where 
\begin{eqnarray}
&&S_{N-1}\left( \hat{\jmath}\right) =\sum_{\substack{ 1\leq
j_{1}<j_{2}<\cdots <j_{N-1}\leq s  \\ j_{1},j_{2},\cdots ,j_{N-1}\neq j}}%
n_{j_{1}}\cdots n_{_{j_{N-1}}} \\
&&\left. {}\right. \quad h_{jk}=-\int \eta _{j}^{\ast }\left( \mathbf{r,\xi }%
\right) \left\{ \sum_{_{1\leq j\leq N_{nuc}}}\frac{Z_{j}}{\left\| \mathbf{R}%
_{j}-\mathbf{r}\right\| }+\frac{1}{2m_{e}}\nabla _{\mathbf{r}}^{2}\right\}
\eta _{k}\left( \mathbf{r,\xi }\right) d\mathbf{r}d\xi ;1\leq j,k\leq 2s \\
&&\left\langle jk||lm\right\rangle =\int \left\{ \eta _{j}^{\ast }\left( 
\mathbf{r}_{1}\mathbf{,\xi }_{1}\right) \eta _{k}^{\ast }\left( \mathbf{r}%
_{2}\mathbf{,\xi }_{2}\right) \frac{1}{\left\| \mathbf{r}_{1}-\mathbf{r}%
_{2}\right\| }\eta _{l}\left( \mathbf{r}_{1}\mathbf{,\xi }_{1}\right) \eta
_{m}\left( \mathbf{r}_{2}\mathbf{,\xi }_{2}\right) \right.  \notag \\
&&\qquad \qquad \qquad -\left. \eta _{j}^{\ast }\left( \mathbf{r}_{1}\mathbf{%
,\xi }_{1}\right) \eta _{k}^{\ast }\left( \mathbf{r}_{2}\mathbf{,\xi }%
_{2}\right) \frac{1}{\left\| \mathbf{r}_{1}-\mathbf{r}_{2}\right\| }\eta
_{l}\left( \mathbf{r}_{2}\mathbf{,\xi }_{1}\right) \eta _{m}\left( \mathbf{r}%
_{1}\mathbf{,\xi }_{2}\right) \right\} d\mathbf{r}_{1}d\mathbf{r}_{2}d\xi
_{1}d\xi _{2}  \notag \\
&&\left. {}\right. \qquad \qquad \qquad \qquad \qquad \qquad \qquad \qquad
\qquad \qquad \quad \qquad 1\leq j,k\leq 2s;1\leq l,m\leq 2s \\
&&\left\langle jk|||jk\right\rangle =\left\langle jk||jk\right\rangle
+\left\langle j\,k+s||j\,k+s\right\rangle +\left\langle
j+s\,k||j+s\,k\right\rangle +\left\langle j+s\,k+s||j+s\,k+s\right\rangle 
\notag \\
&&\left. {}\right. \qquad \qquad \qquad \qquad \qquad \qquad \qquad \qquad
\qquad \qquad \qquad \qquad \qquad \qquad \qquad 1\leq j,k\leq s
\end{eqnarray}%
The total energy, $\mathcal{E}\left( g^{N}\right) $, is then given in the
canonical coordinates as 
\begin{equation}
\mathcal{E}\left( g^{N}\right) =\mathcal{E}\left( \mathbf{\zeta }^{\ast },%
\mathbf{\zeta }\right) =\frac{\mathcal{H}\left( \mathbf{\zeta }^{\ast },%
\mathbf{\zeta }\right) }{\mathcal{N}\left( \mathbf{\zeta }^{\ast },\mathbf{%
\zeta }\right) }.
\end{equation}%
One should observe that this coordinate system \emph{is geminal} \emph{%
dependent }i.e. $\mathbf{\zeta =\zeta }\left( g\right) $ unlike the $\mathbf{%
c}$ and $\mathbf{\gamma }$ coordinate systems and the preceding expressions 
\emph{do not} provide a global representation of the energy function over
the non linear manifold of AGP\ states. However they \emph{do }provide a
procedure for the calculation of the energy of the N-electron AGP state and
its derivatives at a specific point in this manifold.

\subsection{State Energy in Atomic Orbital Basis}

The energy expression eq. $\left( \ref{canen}\right) $ can be easily
transformed to a mixed form in terms of \emph{fixed} atomic orbitals \emph{%
and }canonical coordinates as (see Appendix \ref{Trans2AObasis}) 
\begin{eqnarray}
&&\mathcal{H}\left( \mathbf{\zeta }^{\ast },\mathbf{\zeta }\right) =h_{0}%
\mathcal{N}\left( g^{N}\right)  \notag \\
&&\qquad +\left( N\right) !^{2}\left\{ \sum_{j_{1,}j_{2}=1}^{r}\sum_{1\leq
j\leq s}n_{j}S_{N-1}\left( \hat{\jmath}\right) \left\{ V_{f\,j_{1}j}^{\ast
}V_{f\,j_{2}j}+V_{f\,j_{1}j+s}^{\ast }V_{f\,j_{2}j+s}\right\}
h_{j_{1}j_{2}}^{a}\right.  \notag \\
&&\qquad \qquad +\sum_{j_{1,}j_{2,}j_{3,}j_{4}=1}^{r}\left\{ \overset{}{%
\underset{}{}}\sum_{1\leq j\leq s}n_{j}S_{N-1}\left( \hat{\jmath}\right)
V_{f\,j_{1}j}^{\ast }V_{f\,j_{2}j+s}^{\ast
}V_{f\,j_{3}j}V_{f\,j_{4}j+s}\right.  \notag \\
&&\qquad \qquad \qquad \qquad +\sum_{1\leq j\neq k\leq s}\zeta _{j}^{\ast
}\zeta _{k}S_{N-1}\left( \hat{\jmath}\hat{k}\right) V_{f\,j_{1}j}^{\ast
}V_{f\,j_{2}j+s}^{\ast }V_{f\,j_{3}k}V_{f\,j_{4}k+s}  \notag \\
&&\left. {}\right. \qquad \qquad \quad \left. \left. +\sum_{1\leq j,k\leq
s}n_{j}n_{k}S_{N-2}\left( \hat{\jmath}\hat{k}\right) V_{f\,j_{1}j}^{\ast
}V_{f\,j_{2}k}^{\ast }V_{f\,j_{3}j}V_{f\,j_{4}k}\right\} \left\langle
j_{1}j_{2}||j_{3}j_{4}\right\rangle ^{a}\right\}  \label{AGPen}
\end{eqnarray}%
where the matrix elements $\left\{ h_{j_{1}j_{2}}^{a},\left\langle
j_{1}j_{2}||j_{3}j_{4}\right\rangle ^{a}\right\} $ of the Hamiltonian are
computed in the atomic orbital basis and spin can be integrated out as
displayed in Appendix \ref{SpinInteg}. Hence one can calculate the
unnormalized energy of a $2N$-particle AGP state characterized by geminal
coefficients $\left\{ \gamma _{ij};1\leq i<j\leq r\right\} $ by solving the $%
r\times r$ generalized eigenvalue equation eq. $\left( \ref{NGSO}\right) $
to obtain $\left\{ n_{j};1\leq j\leq s\right\} $ and $\mathbb{V}$, then
using eq. $\left( \ref{eta}\right) $ to produce $\left\{ \zeta _{j},1\leq
j\leq s\right\} $, and finally substituting these values into eq. $\left( %
\ref{AGPen}\right) $ for \emph{fixed values }of the one and two particle
integrals\emph{\ }calculated\emph{\ }in the\emph{\ atomic orbital }basis%
\emph{.}

\section{First Derivatives of the AGP\ State Energy\label{First Derivatives}}

The energy at any point can be expanded in a Taylors series expansion (see
Appendix \ref{DiffProp}) in terms of derivatives with respect to any of the
coordinate systems we have discussed, which can then be transformed between
each other. In particular the derivatives of the AGP energy, $\mathcal{E}%
\left( g^{N}\right) $, in terms of the $\mathbf{\gamma }$ coordinate system
evaluated at a specific geminal $\left| g\right\rangle $ can be expressed in
terms of the derivatives of $\mathcal{E}\left( g^{N}\right) $ in the $%
\mathbf{c}$ coordinate system using the coordinate transformation of eqs. $%
\left( \ref{coordtrans}\right) $ and $\left( \ref{d1trans}\right) $. When
evaluating these derivatives at a particular point $\left| g\right\rangle $
it is convenient to \emph{choose} the coordinate system $\mathbf{c}$ \emph{%
to be }the canonical coordinate system $\mathbf{\zeta }\left( g\right) $ for
the geminal $\left| g\right\rangle $. It should be emphasized that this
approach is \emph{not} the same as differentiating the canonical energy
expression based on eq. $\left( \ref{canen}\right) $ with respect to
canonical coordinates, as the canonical coordinate system is geminal
dependent.

The derivatives of $\mathcal{E}\left( \mathbf{\gamma }^{\ast },\mathbf{%
\gamma }\right) $, which are expressed in terms of the derivatives of $%
\mathcal{H}\left( \mathbf{\gamma }^{\ast },\mathbf{\gamma }\right) $ and $%
\mathcal{N}\left( \mathbf{\gamma }^{\ast },\mathbf{\gamma }\right) $, can be
easily displayed using the second quantized operators introduced in eq. $%
\left( \ref{quantgem}\right) $ and utilizing the relationship 
\begin{equation}
\frac{\partial G\left( \mathbf{\gamma }\right) ^{\dagger N}}{\partial \gamma
_{pq}}=Na_{p}^{\dagger }a_{q}^{\dagger }G\left( \mathbf{\gamma }\right)
^{\dagger N-1}=Na_{p}^{\dagger }a_{q}^{\dagger }G^{\dagger N-1}
\label{GPderiv}
\end{equation}%
Note that in the last term of eq. $\left( \ref{GPderiv}\right) $ $G^{\dagger
N-1}$ is coordinate free and could be expanded in \emph{any }coordinate
system. The derivative of the unnormalized energy in the coordinate system $%
\mathbf{\gamma }$ is thus 
\begin{eqnarray}
\frac{\partial \mathcal{H}\left( g\right) }{\partial \gamma _{pq}} &=&\frac{%
\partial \left\langle G^{\dagger N}\phi |HG^{\dagger N}\phi \right\rangle }{%
\partial \gamma _{pq}}=N\left\langle G^{\dagger N}\phi |Ha_{p}^{\dagger
}a_{q}^{\dagger }G^{\dagger N-1}\phi \right\rangle  \notag \\
\frac{\partial \mathcal{H}\left( g\right) }{\partial \gamma _{pq}^{\ast }}
&=&\frac{\partial \left\langle G^{\dagger N}\phi |HG^{\dagger N}\phi
\right\rangle }{\partial \gamma _{pq}^{\ast }}=N\left\langle a_{p}^{\dagger
}a_{q}^{\dagger }G^{\dagger N-1}\phi |HG^{\dagger N}\phi \right\rangle
\end{eqnarray}%
and are related to the derivatives in the $\mathbf{c}$ coordinate systems by 
\begin{eqnarray}
\frac{\partial \mathcal{H}\left( g\right) }{\partial \mathbf{\gamma }}
&=&\left( \mathbb{W}_{f}\otimes \mathbb{W}_{f}\right) \frac{\partial 
\mathcal{H}\left( g\right) }{\partial \mathbf{c}}  \notag \\
\frac{\partial \mathcal{H}\left( g\right) }{\partial \mathbf{\gamma }^{\ast }%
} &=&\left( \mathbb{W}_{f}^{\ast }\otimes \mathbb{W}_{f}^{\ast }\right) 
\frac{\partial \mathcal{H}\left( g\right) }{\partial \mathbf{c}^{\ast }}
\label{gamder}
\end{eqnarray}%
where $\mathbb{W}_{f}$ is the coordinate transformation from the $\mathbf{%
\gamma }$ to the coordinates $\mathbf{c}$. We now make the important
observation that we can choose $\mathbf{c}$ to be the canonical coordinates
of $\left| g\right\rangle $, so that $\mathbb{W}_{f}=\mathbb{V}_{f}$, and
obtain 
\begin{eqnarray}
\frac{\partial \mathcal{H}\left( g\right) }{\partial \mathbf{c}}
&=&N\left\langle G^{\dagger N}\phi |Hb_{p}^{\dagger }b_{q}^{\dagger
}G^{\dagger N-1}\phi \right\rangle  \notag \\
\frac{\partial \mathcal{H}\left( g\right) }{\partial \mathbf{c}^{\ast }}
&=&N\left\langle b_{p}^{\dagger }b_{q}^{\dagger }G^{\dagger N-1}\phi
|HG^{\dagger N}\phi \right\rangle  \label{cander}
\end{eqnarray}%
where $\left\{ b_{j}^{\dagger };1\leq j\leq r\right\} $ are the creators
associated with the canonical basis $\left\{ \eta _{j};1\leq j\leq r\right\} 
$ of $\left| g\right\rangle $ and hence the right hand side of eq. $\left( %
\ref{cander}\right) $ can be computed in this basis. The derivatives $%
\left\{ \frac{\partial \mathcal{H}\left( g\right) }{\partial \mathbf{\gamma }%
},\frac{\partial \mathcal{H}\left( g\right) }{\partial \mathbf{\gamma }%
^{\ast }}\right\} $ with respect to the geminal coefficients in the atomic
orbital basis can now be obtained by using the transformation displayed in
eq. $\left( \ref{gamder}\right) $ from the canonical to the atomic orbital
basis.

\subsection{Calculation of the Matrix Elements $\left\langle G^{\dagger N}%
\protect\phi |Hb_{p}^{\dagger }b_{q}^{\dagger }G^{\dagger N-1}\protect\phi %
\right\rangle $}

Using the canonical expansion eq. $\left( \ref{canexp}\right) $ one can
develop this matrix element as 
\begin{eqnarray}
&&\left\langle G^{\dagger N}\phi |Hb_{p}^{\dagger }b_{q}^{\dagger
}G^{\dagger N-1}\phi \right\rangle =\left\langle G^{\dagger N}\phi
|Hb_{p}^{\dagger }b_{q}^{\dagger }G^{\dagger N-1}\left( \hat{p}\hat{q}%
\right) \phi \right\rangle =N!\left( N-1\right) !\times  \notag \\
&&\qquad \qquad \sum \zeta _{j_{1}}^{\ast }\cdots \zeta _{j_{N}}^{\ast
}\zeta _{k_{2}}\cdots \zeta _{k_{N}}\langle \eta _{j_{1}}\eta
_{j_{1}+s}\cdots \eta _{j_{N}\,}\eta _{j_{N}+s}\left| H\eta _{p}\eta
_{q}\eta _{k_{2}}\eta _{k_{2}+s}\cdots \eta _{k_{N}\,}\eta
_{k_{N}+s}\right\rangle  \notag \\
&&\overset{1\leq j_{1}<\cdots <j_{N}\leq s;\,1\leq k_{2}<\cdots <k_{N}\leq
s;\,\,k_{i}\neq p,q;\,\,2\leq i\leq N}{}  \label{Derexp}
\end{eqnarray}%
where the carats indicate that terms involving $p$ and $q$ are missing from
the geminal. We introduce the notation $\bar{p}$ to denote the conjugate,
paired GSO\ of $p$ i.e. $\bar{p}=p\pm s$ and $\left| p\right| $ to denote
the index of $p$ defined as 
\begin{equation}
\left| p\right| =\left\{ 
\begin{array}{c}
p\quad \text{if \ }1\leq p\leq s\qquad \\ 
p-s\text{ \ if \ }s+1\leq p\leq 2s%
\end{array}%
\right.
\end{equation}%
We can consider two cases:

\begin{enumerate}
\item $q=\left\vert p\right\vert $
\end{enumerate}

This case can be obtained by differentiating the canonical energy expression
eq. $\left( \ref{canen}\right) $, which we use as a generating function,
with respect to $\zeta _{p}$ to produce 
\begin{eqnarray}
&&\left\langle G^{\dagger N}\phi |Hb_{p}^{\dagger }b_{p+s}^{\dagger
}G^{\dagger N-1}\phi \right\rangle =  \notag \\
&&N!\left( N-1\right) !\left\{ \overset{}{\underset{}{}}S_{N-1}\left( \hat{p}%
\right) \zeta _{p}^{\ast }\left\{ h_{0}+h_{pp}+h_{p+sp+s}+\left\langle
pp+s||pp+s\right\rangle \right\} \right.  \notag \\
&&\left. {}\right. \qquad +\sum_{j\neq p}^{s}S_{N-2}\left( \hat{p}\hat{\jmath%
}\right) \left\{ \overset{}{\underset{}{}}\zeta _{p}^{\ast }n_{j}\left\{
h_{jj}+h_{j+sj+s}+\left\langle jj+s||jj+s\right\rangle +\left\langle
jp|||jp\right\rangle \right\} \right.  \notag \\
&&\left. {}\right. \qquad \qquad \left. +\zeta _{j}^{\ast }\left\langle
jj+s||pp+s\right\rangle \right\} +\sum_{j\neq k\neq p}S_{N-2}\left( \hat{p}%
\hat{\jmath}\hat{k}\right) \zeta _{p}^{\ast }\zeta _{j}^{\ast }\zeta
_{k}\left\langle jj+s||kk+s\right\rangle  \notag \\
&&\qquad \qquad \qquad \qquad \qquad \qquad \qquad \qquad \qquad \left.
+\zeta _{p}^{\ast }\sum_{j<k\neq p}S_{N-3}\left( \hat{p}\hat{\jmath}\hat{k}%
\right) n_{j}n_{k}\left\langle jk|||jk\right\rangle \right\}  \label{pps}
\end{eqnarray}%
or from the expansion eq. $\left( \ref{Derexp}\right) $ and evaluating
matrix elements between Independent Particle States (IPS). (See Appendix $%
\ref{CompFiDer}$ for details)

\begin{enumerate}
\item[2] $q\neq \left\vert p\right\vert $
\end{enumerate}

This case can only be obtained by evaluating matrix elements between IPS's
(See Appendix $\ref{CompFiDer}$ for details). 
\begin{eqnarray}
&&\left\langle G^{\dagger N}\phi |Hb_{p}^{\dagger }b_{q}^{\dagger
}G^{\dagger N-1}\phi \right\rangle =  \notag \\
&&N!\left( N-1\right) !\left\{ \sum_{\kappa \neq \left| p\right| \neq \left|
q\right| }^{s}\left\{ S_{N-1}\left( \hat{\kappa}\hat{p}\hat{q}\right) \zeta
_{\kappa }^{\ast }\left\langle \kappa \kappa +s||pq\right\rangle
+S_{N-2}\left( \hat{\kappa}\hat{p}\hat{q}\right) \left\langle \bar{p}\bar{q}%
||\kappa \kappa +s\right\rangle \zeta _{p}^{\ast }\zeta _{q}^{\ast }\zeta
_{\kappa }\right. \right.  \notag \\
&&\quad \quad \quad \quad \quad \quad \quad \quad \left. +\left\langle \bar{q%
}\kappa ||p\kappa \right\rangle n_{\kappa }\zeta _{q}^{\ast }+\left\langle 
\bar{p}\kappa ||q\kappa \right\rangle n_{\kappa }\zeta _{p}^{\ast }\right\}
+\zeta _{p}^{\ast }\left\langle p\bar{p}||pq\right\rangle S_{N-1}\left( \hat{%
p}\hat{q}\right)  \notag \\
&&\quad \quad \quad \quad \quad \quad \quad \quad \quad \qquad \,\left.
+\zeta _{q}^{\ast }\left\langle q\bar{q}||pq\right\rangle S_{N-1}\left( \hat{%
p}\hat{q}\right) +\left\{ \zeta _{q}^{\ast }h_{\bar{q}p}+\zeta _{p}^{\ast
}h_{\bar{p}q}\right\} S_{N-1}\left( \hat{p}\hat{q}\right) \overset{\underset{%
}{}}{\underset{}{}}\right\}  \label{pq}
\end{eqnarray}

\subsection{Transformation of Derivatives to Atomic Orbital Basis}

For details see Appendix \ref{Trans2AObasis}. In the atomic orbital case we
only have derivatives of the form%
\begin{equation}
\left\langle G^{\dagger N}\phi |Ha_{j_{1}}^{\dagger }a_{j_{2}}^{\dagger
}G^{\dagger N-1}\phi \right\rangle ;1\leq j_{1}<j_{2}\leq r  \label{deriv}
\end{equation}%
which are given in terms of the derivatives expressed in the canonical basis
as%
\begin{equation}
\left\langle G^{\dagger N}\phi |Ha_{j_{1}}^{\dagger }a_{j_{2}}^{\dagger
}G^{\dagger N-1}\phi \right\rangle
=\sum_{p,q=1}^{r}V_{f\,j_{1}p}V_{f\,j_{2}q}\left\langle G^{\dagger N}\phi
|Hb_{p}^{\dagger }b_{q}^{\dagger }G^{\dagger N-1}\phi \right\rangle
\label{ftrm}
\end{equation}%
a sum that goes over the two types of matrix elements, paired terms where $%
p=\left\vert q\right\vert $ computed in eq. $\left( \ref{pq}\right) $ and
unpaired terms $p\neq \left\vert q\right\vert $ computed in eq. $\left( \ref%
{pps}\right) $. The paired term gives 
\begin{eqnarray*}
&&\left\langle G^{\dagger N}\phi |Hb_{p}^{\dagger }b_{p+s}^{\dagger
}G^{\dagger N-1}\phi \right\rangle =N!\left( N-1\right) !\times \\
&&\left\{ h_{0}S_{N-1}\left( \hat{p}\right) \zeta _{p}^{\ast
}+\sum_{j_{1,}j_{2}=1}^{r}\left\{ \underset{}{\overset{}{}}S_{N-1}\left( 
\hat{p}\right) \zeta _{p}^{\ast }\left\{ V_{f\,j_{1}p}^{\ast
}V_{f\,j_{2}p}+V_{f\,j_{1}p+s}^{\ast }V_{f\,j_{2}p+s}\right\}
h_{j_{1}j_{2}}^{a}\right. \right. \\
&&\qquad \qquad \qquad \qquad \qquad \quad \left. +\right. \left.
\sum_{j\neq p}^{s}S_{N-2}\left( \hat{p}\hat{\jmath}\right) \zeta _{p}^{\ast
}n_{j}\left\{ V_{f\,j_{1}j}^{\ast }V_{f\,j_{2}j}+V_{f\,j_{1}j+s}^{\ast
}V_{f\,j_{2}j+s}\right\} h_{j_{1}j_{2}}^{a}\right\}
\end{eqnarray*}%
\begin{eqnarray*}
&&\left. +\right. \sum_{j_{1,}j_{2,}j_{3,}j_{4}=1}^{r}\left\{ \underset{}{%
\overset{}{\left. {}\right. }}S_{N-1}\left( \hat{p}\right) \zeta _{p}^{\ast
}V_{f\,j_{1}p}^{\ast }V_{f\,j_{2}p+s}^{\ast
}V_{f\,j_{3}p}V_{f\,j_{4}p+s}\right. \\
&&\qquad +\sum_{j\neq p}^{s}S_{N-2}\left( \hat{p}\hat{\jmath}\right) \left\{
\zeta _{p}^{\ast }n_{j}\left\{ V_{f\,j_{1}j}^{\ast }V_{f\,j_{2}j+s}^{\ast
}V_{f\,j_{3}j}V_{f\,j_{4}j+s}+V_{f\,j_{1}j}^{\ast }V_{f\,j_{2}p}^{\ast
}V_{f\,j_{3}p}V_{f\,j_{4}j}\right. \right. \\
&&\qquad \qquad \qquad +V_{f\,j_{1}j}^{\ast }V_{f\,j_{2}p+s}^{\ast
}V_{f\,j_{3}p+s}V_{f\,j_{4}j}+V_{f\,j_{1}j+s}^{\ast }V_{f\,j_{2}p}^{\ast
}V_{f\,j_{3}p}V_{f\,j_{4}j+s} \\
&&\qquad \qquad \qquad \qquad \left. +V_{f\,j_{1}j+s}^{\ast
}V_{f\,j_{2}p+s}^{\ast }V_{f\,j_{3}p+s}V_{f\,j_{4}j+s}\right\} \left.
+\,\zeta _{j}^{\ast }V_{f\,j_{1}j}^{\ast }V_{f\,j_{2}j+s}^{\ast
}V_{f\,j_{3}p}V_{f\,j_{4}p+s}\underset{}{\overset{}{}}\right\}
\end{eqnarray*}%
\begin{eqnarray}
&&+\sum_{j\neq k\neq p}^{s}S_{N-2}\left( \hat{p}\hat{\jmath}\hat{k}\right)
\zeta _{p}^{\ast }\zeta _{j}^{\ast }\zeta _{k}V_{f\,j_{1}j}^{\ast
}V_{f\,j_{2}j+s}^{\ast }V_{f\,j_{3}k}V_{f\,j_{4}k+s}\qquad \qquad \qquad 
\notag \\
&&\qquad \quad +\frac{1}{2}\zeta _{p}^{\ast }\sum_{j,k\neq
p}^{s}S_{N-3}\left( \hat{p}\hat{\jmath}\hat{k}\right) n_{j}n_{k}\left\{
V_{f\,j_{1}j}^{\ast }V_{f\,j_{2}k}^{\ast }V_{f\,j_{3}k}V_{f\,j_{4}pj}\right.
\notag \\
&&\qquad \qquad \qquad +V_{f\,j_{1}j}^{\ast }V_{f\,j_{2}k+s}^{\ast
}V_{f\,j_{3}k+s}V_{f\,j_{4}pj}+V_{f\,j_{1}j+s}^{\ast }V_{f\,j_{2}k}^{\ast
}V_{f\,j_{3}k}V_{f\,j_{4}pj+s}  \notag \\
&&\qquad \qquad \qquad \qquad \qquad \qquad \left. \left. \left. \left.
+V_{f\,j_{1}j+s}^{\ast }V_{f\,j_{2}k+s}^{\ast
}V_{f\,j_{3}k+s}V_{f\,j_{4}pj+s}\right\} \underset{}{\overset{}{}}\right\}
\right\} \left\langle j_{1}j_{2}||j_{3}j_{4}\right\rangle ^{a}\underset{}{%
\overset{\underset{}{\overset{}{}}}{}}\right\}
\end{eqnarray}%
and the unpaired term gives 
\begin{eqnarray}
&&\left\langle G^{\dagger N}\phi |Hb_{p}^{\dagger }b_{q}^{\dagger
}G^{\dagger N-1}\phi \right\rangle =N!\left( N-1\right) !\times  \notag \\
&&\left\{ S_{N-1}\left( \hat{p}\hat{q}\right)
\sum_{j_{1},j_{2}=1}^{r}\left\{ V_{f\,j_{1}\bar{q}}^{\ast
}V_{f\,j_{2}p}\zeta _{q}^{\ast }+V_{f\,j_{1}\bar{p}}^{\ast
}V_{f\,j_{2}q}\zeta _{p}^{\ast }\right\} h_{j_{1}j_{2}}^{a}\right.  \notag \\
&&+\sum_{j_{1,}j_{2,}j_{3,}j_{4}=1}^{r}\left\{ \sum_{\kappa \neq \left\vert
p\right\vert \neq \left\vert q\right\vert }^{s}S_{N-1}\left( \hat{\kappa}%
\hat{p}\hat{q}\right) \left\{ V_{f\,j_{1}\kappa }^{\ast }V_{f\,j_{2}\kappa
+s}^{\ast }V_{f\,j_{3}p}V_{f\,j_{4}q}\zeta _{\kappa }^{\ast }\right. \right.
\notag \\
&&\qquad \qquad \qquad \qquad +V_{f\,j_{1}\bar{p}}^{\ast }V_{f\,j_{2}\bar{q}%
}^{\ast }V_{f\,j_{3}\kappa }V_{f\,j_{4}\kappa +s}\zeta _{p}^{\ast }\zeta
_{q}^{\ast }\zeta _{\kappa }+V_{f\,j_{1}\bar{q}}^{\ast }V_{f\,j_{2}\kappa
}^{\ast }V_{f\,j_{3}p}V_{f\,j_{4}\kappa }n_{\kappa }\zeta _{q}^{\ast } 
\notag \\
&&\qquad \qquad \qquad \qquad \qquad \qquad \qquad \qquad \qquad \qquad
\qquad \qquad \left. +V_{f\,j_{1}\bar{p}}^{\ast }V_{f\,j_{2}\kappa }^{\ast
}V_{f\,j_{3}q}V_{f\,j_{4}\kappa }n_{\kappa }\zeta _{p}^{\ast }\right\} 
\notag \\
&&+S_{N-1}\left( \hat{p}\hat{q}\right) \left\{ \zeta _{p}^{\ast
}V_{f\,j_{1}p}^{\ast }V_{f\,j_{2}p+s}^{\ast
}V_{f\,j_{3}p}V_{f\,j_{4}q}+\zeta _{q}^{\ast }V_{f\,j_{1}q}^{\ast
}V_{f\,j_{2}q+s}^{\ast }V_{f\,j_{3}p}V_{f\,j_{4}q}\right\} \left. \left. 
\underset{}{\overset{}{}}\right\} \left\langle
j_{1}j_{2}||j_{3}j_{4}\right\rangle ^{a}\underset{}{\overset{\underset{}{%
\overset{}{}}}{}}\right\}
\end{eqnarray}

\section{Method of Optimization\label{Methods of Optimization}}

In figure 1 we show a schematic of the proposed optimization procedure which
has the following steps:

\begin{figure}[h]
\includegraphics[width=10cm]{agpfig.eps}
\caption{Optimization Flow Diagram}
\label{flowdiag}
\end{figure}

\begin{enumerate}
\item The integrals in the atomic orbital basis are calculated and stored.

\item An initial choice of geminal coefficients in a localized atomic
orbital basis, based on qualitative physical and chemical properties of the
system is made.

\item Using these coefficients, the fixed overlap matrix of the atomic
orbital basis functions and the method described in section \ref{The AGP
State}, values for the canonical coordinates are determined, the number of
computations in this step scales as $r^{3}$.

\item Using the method outlined in \cite{Jensen80, Jensen81} and the
eigenvalues of $\left| g\right\rangle \left\langle g\right| $, produced in
the previous step, the symmetric functions needed for the evaluation of the
energy and the derivatives of the energy with respect to the geminal
coefficients in the atomic orbital basis are calculated, the number of
computations in this step scales as $s^{4}$

\item The derivatives of the energy are now calculated using the
eigenvectors of $\left| g\right\rangle \left\langle g\right| $ produced in
step 3, the integrals produced in step 1 and the quantities $\mathcal{D}%
_{0i_{1}i_{2}},\mathcal{D}_{1i_{1}i_{2}}$ and $\mathcal{D}_{2i_{1}i_{2}}$
described in Appendix \ref{CostofComp}. The number of computations needed to
produce these $\mathcal{D}$ quantities scales as $s^{4}$ and the step
combining these quantities with the integrals scales as $r^{5}$.

\item The gradient of the energy function formed from the partial
derivatives produced in step 5 is then passed to a standard update
optimization procedure which returns a geminal coefficient change vector

\item A new set of geminal coefficients in the atomic orbital basis is
produced and one returns to step 3 and continues the iterative cycle until
convergence is obtained.
\end{enumerate}

\section{Discussion\label{Discussion}}

Introducing correlation corrections to electronic state vectors of molecular
systems is notoriously expensive in terms of computer resources and almost
always involves some type of subjective decision on which terms to keep and
which to neglect. In comparison the choice of an Independent Particle (IP)\
state\ vector is free from any bias, save for the choice of atomic orbital
basis set, is dramatically less expensive to produce and the cost of
computation scales at worst as $r^{4}$, where $r$ is the size of the atomic
spin orbital basis set. In this paper we have described how a correlated
state, the Antisymmetrized Geminal Power (AGP) state is also free of
subjective bias, once the basis set has been determined, and can be chosen
by a computational process that scales as $r^{5}$, which is overall superior
to the scaling of large CASSCF calculations. In an analogous fashion to IP\
state vectors, such an unbiasedly chosen AGP\ state vector can be used as a 
\emph{correlated }reference state from which to introduce further
correlation effects. As the reference would now be a ''dressed state''
perturbation expansions will converge much faster than in the IP case and
lower orders of corrections will provide better approximations to the exact
state vector.

Although the choice of AGP\ state is unbiased, chemical insight can still
play an important role in decreasing the cost of the calculations through
the choice of atomic orbital basis sets. As the optimization procedure
described in this work takes place in a \emph{fixed }atomic orbital basis, a
chemically sensible choice of a localized basis can decrease dramatically
the number of integrals and geminal coefficients that are numerically
significant and thus the integral and geminal coefficient matrices will have
block diagonal form. This will have a similar effect to that of point group
symmetry decreasing the effective size of a basis.

The formalism in this paper utilizes the full generality of a general spin
orbital description which allows the treatment of relativistic and magnetic
effects, and produces a more flexible approximation scheme than spin
restricted ones. It is straightforward, however, to modify the equations so
that they describe spin symmetry constraints.

\section{Acknowledgements}

We would like thank the support of the Petroleum Research Foundation through
a grant PRF 34859-AC6.

\appendix 

\section{Geminal Coordinate Transformations\label{appengct}}

A geminal $\left| g\right\rangle $ can be expanded in the bases $\left\{
\left| \psi _{j}\psi _{k}\right\rangle ;1\leq j<k\leq r\right\} $ and $%
\left\{ \left| \varphi _{j}\varphi _{k}\right\rangle ;1\leq j<k\leq
r\right\} $ as

\begin{equation}
\left| g\right\rangle =\frac{1}{2}\sum_{1\leq j,k\leq r}c_{jk}\left| \psi
_{j}\psi _{k}\right\rangle =\frac{1}{2}\sum_{1\leq j,k\leq r}\gamma
_{jk}\left| \varphi _{j}\varphi _{k}\right\rangle
\end{equation}%
where we have used the antisymmetry of the basis elements and set $%
c_{kj}=-c_{jk}$ and $\gamma _{kj}=-\gamma _{jk}$ for $1\leq j<k\leq r$. In
array of basis function notation this becomes 
\begin{equation}
\left| g\right\rangle =\left( \left| \underline{\psi }\underline{\psi }%
\right\rangle \right) \mathbf{c}=\left( \left| \underline{\varphi }%
\underline{\varphi }\right\rangle \right) \mathbf{\gamma }
\end{equation}%
Hence 
\begin{equation}
\left( \left| \underline{\psi }\underline{\psi }\right\rangle \right) 
\mathbf{c=}\left( \left| \underline{\varphi }\underline{\varphi }%
\right\rangle \right) \left( \mathbb{W}_{f}\otimes \mathbb{W}_{f}\right) 
\mathbf{c}
\end{equation}%
and thus 
\begin{equation}
\mathbf{\gamma =}\left( \mathbb{W}_{f}\otimes \mathbb{W}_{f}\right) \mathbf{%
c=}\mathbb{W}_{f}\mathbf{c}\mathbb{W}_{f}^{t}
\end{equation}

\section{CI Coefficients of AGP States\label{Pfaff}}

An AGP\ state can be expanded in an unordered linearly dependent bases as 
\begin{equation}
\frac{1}{2^{N}}\sum_{1\leq j_{1},j_{2},\cdots ,j_{2N-1},j_{2N}\leq r}\gamma
_{j_{1}j_{2}}\cdots \gamma _{j_{2N-1}j_{2N}}\left\vert \varphi
_{j_{1}}\varphi _{j_{2}}\cdots \varphi _{j_{2N-1}}\varphi
_{j_{2N}}\right\rangle
\end{equation}%
Using the symmetry property 
\begin{equation}
\left\vert \varphi _{j_{1}}\varphi _{j_{2}}\cdots \varphi _{j_{2N-1}}\varphi
_{j_{2N}}\right\rangle =\left( -1\right) ^{\pi \left( \sigma \right)
}\left\vert \varphi _{j_{\sigma \left( 1\right) }}\varphi _{j_{\sigma \left(
2\right) }}\cdots \varphi _{j_{\sigma \left( 2N-1\right) }}\varphi
_{j_{\sigma \left( 2N\right) }}\right\rangle \qquad \forall \sigma \in S_{2N}
\end{equation}%
where $\pi \left( \sigma \right) $ is the parity of the permutation $\sigma $
that belongs to the symmetric group $S_{2N}$, the coefficients of
configurations that can be rearranged into each other can be collected
together as the coefficient, $Pf\left( \gamma :j_{1}\cdots j_{2N}\right) $,
of a single ordered configuration $\left\vert \varphi _{j_{1}}\cdots \varphi
_{j_{2N}}\right\rangle $ where $j_{1}<j_{2}<\cdots <j_{2N-1}<j_{2N}$. This
coefficient is given by 
\begin{eqnarray}
&&Pf\left( \gamma :j_{1}\cdots j_{2N}\right) \left\vert \varphi
_{j_{1}}\varphi _{j_{2}}\cdots \varphi _{j_{2N-1}}\varphi
_{j_{2N}}\right\rangle  \notag \\
&=&\frac{1}{N!2^{N}}\sum_{\sigma \in S_{2N}}\gamma _{j_{\sigma \left(
1\right) }j_{\sigma \left( 2\right) }}\cdots \gamma _{j_{\sigma \left(
2N-1\right) }j_{\sigma \left( 2N\right) }}\left\vert \varphi _{j_{\sigma
\left( 1\right) }}\varphi _{j_{\sigma \left( 2\right) }}\cdots \varphi
_{j_{\sigma \left( 2N-1\right) }}\varphi _{j_{\sigma \left( 2N\right)
}}\right\rangle  \notag \\
&=&\frac{1}{N!2^{N}}\sum_{\sigma \in S_{2N}}\left( -1\right) ^{\pi \left(
\sigma \right) }\gamma _{j_{\sigma \left( 1\right) }j_{\sigma \left(
2\right) }}\cdots \gamma _{j_{\sigma \left( 2N-1\right) }j_{\sigma \left(
2N\right) }}\left\vert \varphi _{j_{1}}\varphi _{j_{2}}\cdots \varphi
_{j_{2N-1}}\varphi _{j_{2N}}\right\rangle
\end{eqnarray}%
where 
\begin{eqnarray}
Pf\left( \gamma :j_{1}\cdots j_{2N}\right) &=&Pf\left( 
\begin{array}{ccccc}
0 & \gamma _{j_{1}j_{2}} & \gamma _{j_{1}j_{4}} & \cdots & \gamma
_{j_{1}j_{2N}} \\ 
\gamma _{j_{2}j_{1}} & 0 & \gamma _{j_{2}j_{4}} & \vdots & \gamma
_{j_{2}j_{2N}} \\ 
\vdots & \vdots & \ddots & \vdots & \vdots \\ 
\vdots & \vdots & \vdots & 0 & \gamma _{j_{2N-1}j_{2N}} \\ 
\gamma _{j_{2N}j_{1}} & \cdots & \cdots & \gamma _{j_{2N}j_{2N-1}} & 0%
\end{array}%
\right)  \notag \\
&=&\frac{1}{N!2^{N}}\sum_{\sigma \in S_{2N}}\left( -1\right) ^{\pi \left(
\sigma \right) }\gamma _{j_{\sigma \left( 1\right) }j_{\sigma \left(
2\right) }}\cdots \gamma _{j_{\sigma \left( 2N-1\right) }j_{\sigma \left(
2N\right) }}
\end{eqnarray}%
The symbol stands for Pfaffian, which is a function that associates a scalar
with antisymmetric matrices. For further details on Pfaffians see \cite%
{Weiner94,Marcus73,Caianiello73}

\section{Canonical Coordinates Construction\label{Appenccc}}

Bloch and Messiah \cite{Bloch62} and Zumino \cite{Zumino62} proved

\begin{theorem}
If $\mathbf{c}$ is a $2s\times 2s$ complex antisymmetric matrix then there
exists a unitary matrix $\mathbb{T}$ such that 
\begin{equation}
\mathbb{T}^{t}\mathbf{c}\mathbb{T=}\left( 
\begin{array}{cc}
\mathbf{0} & \mathbf{\zeta } \\ 
-\mathbf{\zeta } & \mathbf{0}%
\end{array}%
\right)  \label{tct}
\end{equation}%
where 
\begin{equation}
\mathbf{\zeta =}\left( 
\begin{array}{ccc}
\zeta _{1} & 0 & 0 \\ 
0 & \ddots & 0 \\ 
0 & 0 & \zeta _{s}%
\end{array}%
\right) ;\qquad \zeta _{j}\in \mathbb{C},\quad 1\leq j\leq s
\end{equation}
\end{theorem}

Hence from eq. $\left( \ref{tct}\right) $ 
\begin{equation}
\mathbb{T}^{\dagger }\mathbf{cc}^{\dagger }\mathbb{T}=\left( 
\begin{array}{cc}
\mathbf{n} & \mathbf{0} \\ 
\mathbf{0} & \mathbf{n}%
\end{array}%
\right)  \label{TccT}
\end{equation}%
however, eq. $\left( \ref{TccT}\right) $ does not imply eq. $\left( \ref{tct}%
\right) $, which shows that from the class of unitary transformations that
diagonalizes $\mathbf{cc}^{\dagger }$ only a subset of these satisfy eq. $%
\left( \ref{tct}\right) $. In order to obtain $\mathbb{T}^{\prime }s$ that
do satisfy eq. $\left( \ref{tct}\right) $ they displayed that

\begin{itemize}
\item any unitary transformation that diagonalizes $\mathbf{cc}^{\dagger }$
produces an antisymmetric block diagonal matrix $\mathbb{U}^{t}\mathbf{c}%
\mathbb{U}$, where the size of the blocks are equal to the level of the
degeneracy of the eigenvalues of $\mathbf{cc}^{\dagger }$, and whose real
and imaginary parts \emph{commute, }

\item the blocks have the property 
\begin{equation}
\mathbb{C}_{\kappa }\mathbb{C}_{\kappa }^{\dagger }=\eta _{\kappa }\mathbb{I}%
_{d\left( \kappa \right) };\quad 1\leq \kappa \leq p
\end{equation}%
where $p$ is the number of blocks, $\eta _{\kappa }$ is the value of the $%
d\left( \kappa \right) $-degenerate $\kappa ^{th}$ eigenvalue of $\mathbf{cc}%
^{\dagger }$,

\item $\mathbb{U}^{t}\mathbf{c}\mathbb{U}$ can be brought to skew diagonal
form by a \emph{real orthogonal} transformation $\mathbb{O}^{t}\mathbb{U}^{t}%
\mathbf{c}\mathbb{UO}$, where $\mathbb{O}$ commutes with $\mathbb{U}%
^{\dagger }\mathbf{cc}^{\dagger }\mathbb{U}$ i.e. 
\begin{equation}
\mathbb{O}^{t}\left( 
\begin{array}{cc}
\mathbf{n} & \mathbf{0} \\ 
\mathbf{0} & \mathbf{n}%
\end{array}%
\right) \mathbb{O=}\left( 
\begin{array}{cc}
\mathbf{n} & \mathbf{0} \\ 
\mathbf{0} & \mathbf{n}%
\end{array}%
\right) \mathbb{O}^{t}\mathbb{O=}\left( 
\begin{array}{cc}
\mathbf{n} & \mathbf{0} \\ 
\mathbf{0} & \mathbf{n}%
\end{array}%
\right)
\end{equation}%
and the transformation $\mathbb{O}$ is a simple rearrangement of the direct
sum of the transformations $\mathbb{O}_{\kappa }$ that bring the individual
antisymmetric blocks, $\mathbb{C}_{\kappa }$, of $\mathbb{U}^{t}\mathbf{c}%
\mathbb{U}$ to skew diagonal form and in particular simultaneously
diagonalize the real and imaginary parts of $\mathbb{C}_{\kappa }$ and

\item hence $\mathbb{T=UO}$
\end{itemize}

In order to compute the transformations $\mathbb{O}_{\kappa }$ we consider
the real and imaginary parts of the \emph{antisymmetric unitary matrix} $%
\mathbb{B}_{\kappa }=\eta _{\kappa }^{-\frac{1}{2}}\mathbb{C}_{\kappa }$ and
denote them by%
\begin{eqnarray}
\mathbb{B}_{R} &=&\func{Re}\mathbb{B}_{\kappa }  \notag \\
\mathbb{B}_{I} &=&\func{Im}\mathbb{B}_{\kappa }
\end{eqnarray}%
where we drop the subscript $\kappa $, which plays no role in the
considerations of this section. We know that 
\begin{equation}
\mathbb{B}_{R}=-\mathbb{B}_{R}^{t}\text{ and }\mathbb{B}_{I}=-\mathbb{B}%
_{I}^{t}
\end{equation}%
and 
\begin{eqnarray}
&&\mathbb{BB}^{\ast }=\mathbb{B^{\ast }B}=-\mathbb{I}  \notag \\
&\Rightarrow &\mathbb{B}^{-1}=-\mathbb{B^{\ast }}\text{ and }\mathbb{B}%
_{R}^{2}+\mathbb{B}_{I}^{2}=-\mathbb{I}
\end{eqnarray}%
We can then form the commuting hermitian matrices $\mathbb{H}_{R}$ and $%
\mathbb{H}_{I}$ by 
\begin{equation}
\mathbb{H}_{R}=i\mathbb{B}_{R}\text{ and }\mathbb{H}_{I}=i\mathbb{B}_{I}
\end{equation}%
so that 
\begin{equation}
\mathbb{H}_{R}^{2}+\mathbb{H}_{I}^{2}=\mathbb{I}\text{ and }\left[ \mathbb{H}%
_{R},\mathbb{H}_{I}\right] =0
\end{equation}%
Hence we can conclude that there exists a matrix $\mathbf{\Omega }$ such
that 
\begin{equation}
\mathbb{H}_{R}=\cos \mathbf{\Omega }\text{ and }\mathbb{H}_{I}=\sin \mathbf{%
\Omega }
\end{equation}%
and thus that 
\begin{equation}
\mathbb{H}_{R}\mathbb{H}_{I}=\mathbb{H}_{I}\mathbb{H}_{R}=\sin \mathbf{%
\Omega }\cos \mathbf{\Omega =}\frac{1}{2}\sin 2\mathbf{\Omega }
\end{equation}%
Now 
\begin{equation}
\mathbb{H}_{R}\mathbb{H}_{I}=-\mathbb{B}_{R}\mathbb{B}_{I}
\end{equation}%
so the solution of the \emph{real symmetric} eigenvalue problem 
\begin{equation}
\mathbb{B}_{R}\mathbb{B}_{I}\mathbf{v}_{j}=-\sin 2\mathbf{\Omega \,v}%
_{j}=\lambda _{j}\mathbf{v}_{j}
\end{equation}%
gives a real orthonormal basis $\left\{ \mathbf{v}_{j}\right\} $ in which $%
\sin 2\mathbf{\Omega }$, $\sin \mathbf{\Omega }$, $\cos \mathbf{\Omega }$, $%
\mathbb{H}_{R}$ and $\mathbb{H}_{I}$ are diagonal and $\mathbb{B}_{R}$ and $%
\mathbb{B}_{I}$ are skew diagonal. The eigenvalues $\left\{ \lambda
_{j}\right\} $ are at least doubly degenerate as $\mathbb{B}_{R}$ and $%
\mathbb{B}_{I}$ are antisymmetric. The vectors $\left\{ \mathbf{v}%
_{j}\right\} $ form the columns of $\mathbb{O}$. The diagonal entries
(eigenvalues) of $\mathbb{H}_{R}$ and $\mathbb{H}_{I}$ are given in terms of 
\begin{equation}
\varphi _{j}=\frac{1}{2}\sin ^{-1}\left( -\lambda _{j}\right)
\end{equation}%
as 
\begin{equation}
\alpha _{j}=-i\cos \varphi _{j}=-i\cos \left\{ \frac{1}{2}\sin ^{-1}\left(
-\lambda _{j}\right) \right\}
\end{equation}%
and 
\begin{equation}
\beta _{j}=-i\sin \varphi _{j}=-i\sin \left\{ \frac{1}{2}\sin ^{-1}\left(
-\lambda _{j}\right) \right\}
\end{equation}%
leading to the eigenvalues of $\mathbb{B}$ being 
\begin{equation}
-i\cos \varphi _{j}+\sin \varphi _{j}=-i\left( \cos \varphi _{j}+i\sin
\varphi _{j}\right) =e^{i\left\{ \varphi _{j}+\pi \right\} }=e^{i\theta _{j}}
\end{equation}%
i.e. $\theta _{j}=\varphi _{j}+\pi $.

\section{Energy of the 2N-Electron Unnormalized AGP State\label{EnDetail}}

The Hamiltonian in the canonical basis can be written as 
\begin{equation}
H=h_{0}+\sum_{i,j=1}^{r}h_{ij}b_{i}^{\dagger }b_{j}+\frac{1}{4}%
\sum_{i,j,k,l=1}^{r}\left\langle ij||kl\right\rangle b_{i}^{\dagger
}b_{j}^{\dagger }b_{l}b_{k}
\end{equation}%
giving the 2N-electron unnormalized AGP state energy\ 
\begin{equation}
\left\langle G^{\dagger N}\phi |HG^{\dagger N}\phi \right\rangle
=\left\langle G^{\dagger N}\phi \left| \left\{
h_{0}+\sum_{i,j=1}^{r}h_{ij}b_{i}^{\dagger }b_{j}+\frac{1}{4}%
\sum_{i,j,k,l=1}^{r}\left\langle ij||kl\right\rangle b_{i}^{\dagger
}b_{j}^{\dagger }b_{l}b_{k}\right\} G^{\dagger N}\phi \right. \right\rangle
\end{equation}%
that can be expanded as 
\begin{eqnarray*}
&&\sum_{1\leq j_{1},\cdots ,j_{N}\leq s;\,1\leq k_{1},\cdots ,k_{N}\leq
s}\zeta _{j_{1}}^{\ast }\cdots \zeta _{j_{N}}^{\ast }\zeta _{k_{1}}\cdots
\zeta _{k_{N}} \\
&&\left\langle \eta _{j_{1}}\eta _{j_{1}+s}\cdots \eta _{j_{N}\,}\eta
_{j_{N}+s}\left| \left\{ h_{0}+\sum_{i,j=1}^{r}h_{ij}b_{i}^{\dagger }b_{j}+%
\frac{1}{4}\sum_{i,j,k,l=1}^{r}\left\langle ij||kl\right\rangle
b_{i}^{\dagger }b_{j}^{\dagger }b_{l}b_{k}\right\} \eta _{k_{1}}\eta
_{k_{1}+s}\cdots \eta _{k_{N}\,}\eta _{k_{N}+s}\right. \right\rangle
\end{eqnarray*}%
\begin{equation}
\qquad \qquad \qquad \qquad \qquad \qquad \qquad \qquad \qquad \qquad \qquad
\qquad \qquad
\end{equation}%
The first term is easy to compute--all paired indices must match, and as
they are unordered we pick up a factor $N!$ each from the bra and ket to
obtain 
\begin{eqnarray}
&&h_{0}\sum \zeta _{j_{1}}^{\ast }\cdots \zeta _{j_{N}}^{\ast }\zeta
_{k_{1}}\cdots \zeta _{k_{N}}\left\langle \eta _{j_{1}}\eta _{j_{1}+s}\cdots
\eta _{j_{N}\,}\eta _{j_{N}+s}\right. \left| \eta _{k_{1}}\eta
_{k_{1}+s}\cdots \eta _{k_{N}\,}\eta _{k_{N}+s}\right\rangle =h_{0}\left(
N!\right) ^{2}S_{N}  \notag \\
&&\qquad \qquad \qquad \qquad \qquad \qquad \qquad \qquad \qquad \qquad
\qquad _{\quad \qquad 1\leq j_{1},\cdots ,j_{N}\leq s;\,1\leq k_{1},\cdots
,k_{N}\leq s}
\end{eqnarray}%
The one electron terms are obtained by noting that $b_{i}^{\dagger }b_{j}$
must annihilate the same GSO from the bra and the ket producing 
\begin{eqnarray}
&&\qquad \sum \zeta _{j_{1}}^{\ast }\cdots \zeta _{j_{N}}^{\ast }\zeta
_{k_{1}}\cdots \zeta _{k_{N}}\left\langle \eta _{j_{1}}\eta _{j_{1}+s}\cdots
\eta _{j_{N}\,}\eta _{j_{N}+s}\left| \sum_{i,,j=1}^{r}h_{ij}b_{i}^{\dagger
}b_{j}\eta _{k_{1}}\eta _{k_{1}+s}\cdots \eta _{k_{N}\,}\eta
_{k_{N}+s}\right. \right\rangle  \notag \\
&&_{1\leq j_{1},\cdots ,j_{N}\leq s;\,1\leq k_{1},\cdots ,k_{N}\leq s} 
\notag \\
&&\left. {}\right. \qquad \qquad \qquad \qquad \qquad \qquad \qquad \qquad
\qquad =\left( N!\right) ^{2}\sum_{i=1}^{s}\left\{ h_{ii}+h_{i+si+s}\right\}
n_{i}S_{N-1}\left( \hat{\imath}\right)
\end{eqnarray}%
With the two electron term 
\begin{equation*}
\sum \zeta _{j_{1}}^{\ast }\cdots \zeta _{j_{N}}^{\ast }\zeta _{k_{1}}\cdots
\zeta _{k_{N}}\left\langle \eta _{j_{1}}\eta _{j_{1}+s}\cdots \eta
_{j_{N}\,}\eta _{j_{N}+s}\left| \frac{1}{4}\sum_{i,j,k,l=1}^{r}\left\langle
ij||kl\right\rangle b_{i}^{\dagger }b_{j}^{\dagger }b_{l}b_{k}\eta
_{k_{1}}\eta _{k_{1}+s}\cdots \eta _{k_{N}\,}\eta _{k_{N}+s}\right.
\right\rangle
\end{equation*}%
\begin{equation}
\qquad \qquad \qquad \qquad \qquad \qquad \qquad \qquad \qquad \qquad \qquad
\qquad \qquad \qquad _{1\leq j_{1},\cdots ,j_{N}\leq s;\,1\leq k_{1},\cdots
,k_{N}\leq s\,}
\end{equation}%
we have more choices. To see this clearly we can split the two-electron
operator into the following portions%
\begin{eqnarray}
&&\sum_{i,j,k,l=1}^{r}\left\langle ij||kl\right\rangle b_{i}^{\dagger
}b_{j}^{\dagger }b_{l}b_{k}=  \notag \\
&&\left. {}\right. \sum_{i=1}^{s}\left\langle ii+s||ii+s\right\rangle
\left\{ b_{i}^{\dagger }b_{i+s}^{\dagger }b_{i+s}b_{i}-b_{i}^{\dagger
}b_{i+s}^{\dagger }b_{i}b_{i+s}-b_{i+s}^{\dagger }b_{i}^{\dagger
}b_{i+s}b_{i}+b_{i+s}^{\dagger }b_{i}^{\dagger }b_{i}b_{i+s}\right\}  \notag
\\
&&+\sum_{i\neq l=1}^{s}\left\langle ii+s||ll+s\right\rangle \left\{
b_{i}^{\dagger }b_{i+s}^{\dagger }b_{l+s}b_{l}-b_{i}^{\dagger
}b_{i+s}^{\dagger }b_{l}b_{l+s}-b_{i+s}^{\dagger }b_{i}^{\dagger
}b_{l+s}b_{l}+b_{i+s}^{\dagger }b_{i}^{\dagger }b_{l}b_{l+s}\right\}
\label{twoel1}
\end{eqnarray}%
and 
\begin{eqnarray}
&&+\sum_{i,,j=1}^{s}\left\langle ij||ij\right\rangle \left\{ b_{i}^{\dagger
}b_{j}^{\dagger }b_{j}b_{i}-b_{i}^{\dagger }b_{j}^{\dagger
}b_{i}b_{j}-b_{j}^{\dagger }b_{i}^{\dagger }b_{j}b_{i}+b_{j}^{\dagger
}b_{i}^{\dagger }b_{i}b_{j}\right\}  \notag \\
&&+\sum_{i,j=1}^{s}\left\langle i+sj||i+sj\right\rangle \left\{
b_{i+s}^{\dagger }b_{j}^{\dagger }b_{j}b_{i+s}-b_{i+s}^{\dagger
}b_{j}^{\dagger }b_{i+s}b_{j}-b_{j}^{\dagger }b_{i+s}^{\dagger
}b_{j}b_{i+s}+b_{j}^{\dagger }b_{i+s}^{\dagger }b_{i+s}b_{j}\right\}  \notag
\\
&&+\sum_{i,j}^{s}\left\langle ij+s||ij+s\right\rangle \left\{ b_{i}^{\dagger
}b_{j+s}^{\dagger }b_{j+s}b_{i}-b_{i}^{\dagger }b_{j+s}^{\dagger
}b_{i}b_{j+s}-b_{j+s}^{\dagger }b_{i}^{\dagger
}b_{j+s}b_{i}+b_{j+s}^{\dagger }b_{i}^{\dagger }b_{i}b_{j+s}\right\}  \notag
\\
&&+\sum_{i,j=1}^{s}\left\langle i+sj+s||i+sj+s\right\rangle \left\{
b_{i+s}^{\dagger }b_{j+s}^{\dagger }b_{j+s}b_{i+s}-b_{i+s}^{\dagger
}b_{j+s}^{\dagger }b_{i+s}b_{j+s}\right.  \notag \\
&&\qquad \qquad \qquad \qquad \qquad \qquad \left. -b_{j+s}^{\dagger
}b_{i+s}^{\dagger }b_{j+s}b_{i+s}+b_{j+s}^{\dagger }b_{i+s}^{\dagger
}b_{i+s}b_{j+s}\right\} ;\quad i\neq j\pm s  \notag \\
&&+\text{ other terms}  \label{twoel2}
\end{eqnarray}%
The portion eq. $\left( \ref{twoel1}\right) $ produces 
\begin{equation}
4\left( N\right) !^{2}\left\{ \sum_{1\leq i\leq s}n_{i}S_{N-1}\left( \hat{%
\imath}\right) \left\langle ii+s||ii+s\right\rangle +\sum_{1\leq i\neq j\leq
s}\zeta _{i}^{\ast }\zeta _{j}S_{N-2}\left( \hat{\imath}\hat{\jmath}\right)
\left\langle ii+s||jj+s\right\rangle \right\}
\end{equation}%
and the portion eq. $\left( \ref{twoel2}\right) $ produces%
\begin{equation}
4\left( N\right) !^{2}\sum_{1\leq i,j\leq s}n_{i}n_{j}S_{N-2}\left( \hat{%
\imath}\hat{\jmath}\right) \left\langle ij|||ij\right\rangle
\end{equation}%
and the expectation value of the ''other terms'' is zero. The full
unnormalized energy expression is then 
\begin{eqnarray}
&&\mathcal{H}\left( \mathbf{\zeta }^{\ast },\mathbf{\zeta }\right) =h_{0}%
\mathcal{N}\left( g^{N}\right) +\left( N\right) !^{2}\left\{ \sum_{1\leq
i\leq s}n_{i}S_{N-1}\left( \hat{\imath}\right) \left\{
h_{ii}+h_{i+si+s}+\left\langle ii+s||ii+s\right\rangle \right\} \right. 
\notag \\
&&\qquad \qquad +\sum_{1\leq i\neq j\leq s}\zeta _{i}^{\ast }\zeta
_{j}S_{N-1}\left( \hat{\imath}\hat{\jmath}\right) \left\langle
ii+s||jj+s\right\rangle \left. +\sum_{1\leq i,j\leq
s}n_{i}n_{j}S_{N-2}\left( \hat{\imath}\hat{\jmath}\right) \left\langle
ij|||ij\right\rangle \right\}
\end{eqnarray}

\section{Differential Properties\label{DiffProp}}

\subsection{Taylor Series Expansion}

The functions appearing in eq. $\left( \ref{EnFunc}\right) $ have second
order Taylor series expansions given by 
\begin{equation}
f\left( \mathbf{z}^{\ast }+\mathbf{h}^{\ast },\mathbf{z}+\mathbf{h}\right)
\simeq f\left( \mathbf{z}^{\ast },\mathbf{z}\right) +D^{1}f\left( \mathbf{z}%
^{\ast },\mathbf{z}\right) \bullet \mathbf{h}+\frac{1}{2}\mathbf{h}^{\dagger
}\bullet D^{2}f\left( \mathbf{z}^{\ast },\mathbf{z}\right) \bullet \mathbf{h}
\end{equation}%
where%
\begin{eqnarray}
\mathbf{y} &\mathbf{=}&\left( 
\begin{array}{c}
\mathbf{z}^{\ast } \\ 
\mathbf{z}%
\end{array}%
\right)  \notag \\
D^{1}f\left( \mathbf{z}^{\ast },\mathbf{z}\right) &=&\left( 
\begin{array}{cc}
\frac{\partial f}{\partial \mathbf{z}^{\ast }} & \frac{\partial f}{\partial 
\mathbf{z}}%
\end{array}%
\right)  \notag \\
D^{2}f\left( \mathbf{z}^{\ast },\mathbf{z}\right) &=&\left( 
\begin{array}{cc}
\frac{\partial ^{2}f}{\partial \mathbf{z}^{\ast }\partial \mathbf{z}} & 
\frac{\partial ^{2}f}{\partial \mathbf{z}^{\ast }\partial \mathbf{z}^{\ast }}
\\ 
\frac{\partial ^{2}f}{\partial \mathbf{z}\partial \mathbf{z}} & \frac{%
\partial ^{2}f}{\partial \mathbf{z}\partial \mathbf{z}^{\ast }}%
\end{array}%
\right)
\end{eqnarray}%
\begin{eqnarray}
\left( \frac{\partial f}{\partial \mathbf{z}}\right) _{j} &=&\frac{\partial f%
}{\partial z_{j}}  \notag \\
\left( \frac{\partial ^{2}f}{\partial \mathbf{z}^{\ast }\partial \mathbf{z}}%
\right) _{jk} &=&\frac{\partial ^{2}f}{\partial z_{j}^{\ast }\partial z_{k}}
\end{eqnarray}%
and all the derivatives are evaluated at $\left( \mathbf{z}^{\ast },\mathbf{z%
}\right) $ and all other terms are obtained by complex conjugation.

\subsection{Coordinate Transformation}

The derivatives in different coordinate systems $\mathbf{c}$ and $\mathbf{%
\gamma }$ are related to each by 
\begin{eqnarray}
&&\left( 
\begin{array}{c}
\frac{\partial f}{\partial \mathbf{\gamma }^{\ast }} \\ 
\frac{\partial f}{\partial \mathbf{\gamma }}%
\end{array}%
\right) =\left( 
\begin{array}{cc}
\mathbb{W}_{f}^{\ast }\otimes \mathbb{W}_{f}^{\ast } & \mathbf{0} \\ 
\mathbf{0} & \mathbb{W}_{f}\otimes \mathbb{W}_{f}%
\end{array}%
\right) \left( 
\begin{array}{c}
\frac{\partial f}{\partial \mathbf{c}^{\ast }} \\ 
\frac{\partial f}{\partial \mathbf{c}}%
\end{array}%
\right) =\left( 
\begin{array}{c}
\mathbb{W}_{f}^{\ast }\frac{\partial f}{\partial \mathbf{c}^{\ast }}\mathbb{W%
}_{f}^{\dagger } \\ 
\mathbb{W}_{f}\frac{\partial f}{\partial \mathbf{c}}\mathbb{W}_{f}^{t}%
\end{array}%
\right)  \label{d1trans} \\
&&\left( 
\begin{array}{cc}
\frac{\partial ^{2}f}{\partial \mathbf{\gamma }\partial \mathbf{\gamma }%
^{\ast }} & \frac{\partial ^{2}f}{\partial \mathbf{\gamma }\partial \mathbf{%
\gamma }} \\ 
\frac{\partial ^{2}f}{\partial \mathbf{\gamma }^{\ast }\partial \mathbf{%
\gamma }^{\ast }} & \frac{\partial ^{2}f}{\partial \mathbf{\gamma }^{\ast
}\partial \mathbf{\gamma }}%
\end{array}%
\right) =  \notag \\
&&\left( 
\begin{array}{cc}
\mathbb{W}_{f}\otimes \mathbb{W}_{f} & \mathbf{0} \\ 
\mathbf{0} & \mathbb{W}_{f}^{\ast }\otimes \mathbb{W}_{f}^{\ast }%
\end{array}%
\right) ^{\dagger }\left( 
\begin{array}{cc}
\frac{\partial ^{2}f}{\partial \mathbf{c}\partial \mathbf{c}^{\ast }} & 
\frac{\partial ^{2}f}{\partial \mathbf{c}\partial \mathbf{c}} \\ 
\frac{\partial ^{2}f}{\partial \mathbf{c}^{\ast }\partial \mathbf{c}^{\ast }}
& \frac{\partial ^{2}f}{\partial \mathbf{c}^{\ast }\partial \mathbf{c}}%
\end{array}%
\right) \left( 
\begin{array}{cc}
\mathbb{W}_{f}^{\dagger }\otimes \mathbb{W}_{f}^{\dagger } & \mathbf{0} \\ 
\mathbf{0} & \mathbb{W}_{f}^{t}\otimes \mathbb{W}_{f}^{t}%
\end{array}%
\right) =  \notag \\
&&\qquad \qquad \qquad \qquad \qquad \left( 
\begin{array}{cc}
\mathbb{W}_{f}\otimes \mathbb{W}_{f}\frac{\partial ^{2}f}{\partial \mathbf{c}%
\partial \mathbf{c}^{\ast }}\mathbb{W}_{f}^{\dagger }\otimes \mathbb{W}%
_{f}^{\dagger } & \mathbb{W}_{f}\otimes \mathbb{W}_{f}\frac{\partial ^{2}f}{%
\partial \mathbf{c}\partial \mathbf{c}}\mathbb{W}_{f}^{t}\otimes \mathbb{W}%
_{f}^{t} \\ 
\mathbb{W}_{f}^{\ast }\otimes \mathbb{W}_{f}^{\ast }\frac{\partial ^{2}f}{%
\partial \mathbf{c}^{\ast }\partial \mathbf{c}^{\ast }}\mathbb{W}%
_{f}^{\dagger }\otimes \mathbb{W}_{f}^{\dagger } & \mathbb{W}_{f}^{\ast
}\otimes \mathbb{W}_{f}^{\ast }\frac{\partial ^{2}f}{\partial \mathbf{c}%
^{\ast }\partial \mathbf{c}}\mathbb{W}_{f}^{t}\otimes \mathbb{W}_{f}^{t}%
\end{array}%
\right)  \label{d2trans}
\end{eqnarray}

\subsection{Derivatives of Quotients}

The energy is expressed as a rational function 
\begin{equation}
\mathcal{E}\left( \mathbf{z}^{\ast },\mathbf{z}\right) =\frac{\mathcal{H}%
\left( \mathbf{z}^{\ast },\mathbf{z}\right) }{\mathcal{N}\left( \mathbf{z}%
^{\ast },\mathbf{z}\right) }
\end{equation}%
which has first derivative 
\begin{equation}
D^{1}\mathcal{E}\left( \mathbf{z}^{\ast },\mathbf{z}\right) \mathcal{=}\frac{%
1}{\mathcal{N}\left( \mathbf{z}^{\ast },\mathbf{z}\right) }\left\{ D^{1}%
\mathcal{H}\left( \mathbf{z}^{\ast },\mathbf{z}\right) \mathcal{-E}\left( 
\mathbf{z}^{\ast },\mathbf{z}\right) D^{1}\mathcal{N}\left( \mathbf{z}^{\ast
},\mathbf{z}\right) \right\}
\end{equation}%
and second derivative 
\begin{eqnarray}
D^{2}\mathcal{E}\left( \mathbf{z}^{\ast },\mathbf{z}\right) &=&\frac{1}{%
\mathcal{N}\left( \mathbf{z}^{\ast },\mathbf{z}\right) }\left\{ D^{2}%
\mathcal{H}\left( \mathbf{z}^{\ast },\mathbf{z}\right) \mathcal{-E}\left( 
\mathbf{z}^{\ast },\mathbf{z}\right) D^{2}\mathcal{N}\left( \mathbf{z}^{\ast
},\mathbf{z}\right) \right.  \notag \\
&&\left. -\left\{ D^{1}\mathcal{E}\left( \mathbf{z}^{\ast },\mathbf{z}%
\right) \otimes D^{1}\mathcal{N}\left( \mathbf{z}^{\ast },\mathbf{z}\right)
+D^{1}\mathcal{N}\left( \mathbf{z}^{\ast },\mathbf{z}\right) \otimes D^{1}%
\mathcal{E}\left( \mathbf{z}^{\ast },\mathbf{z}\right) \right\} \right\}
\end{eqnarray}

\section{Computation of First Derivatives\label{CompFiDer}}

\subsection{\textit{Calculation of} $\left\langle G^{\dagger N}\protect\phi %
|Hb_{p}^{\dagger }b_{p+s}^{\dagger }G^{\dagger N-1}\protect\phi %
\right\rangle $}

The 2N-electron AGP unnormalized state energy derivative wrt $\gamma _{pq}$
needs the matrix elements $\left\langle G^{\dagger N}\phi |Hb_{p}^{\dagger
}b_{p+s}^{\dagger }G^{\dagger N-1}\phi \right\rangle $ and $\left\langle
G^{\dagger N}\phi |Hb_{p}^{\dagger }b_{q}^{\dagger }G^{\dagger N-1}\phi
\right\rangle $ for $\left| p\right| \neq \left| q\right| $. We first
consider the $\left( p,p+s\right) $ elements 
\begin{eqnarray*}
&&\left\langle G^{\dagger N}\phi |Hb_{p}^{\dagger }b_{p+s}^{\dagger
}G^{\dagger N-1}\phi \right\rangle \quad \qquad \qquad \qquad \qquad \qquad
\qquad \qquad \qquad \qquad \qquad \qquad \qquad \quad \\
&=&\left\langle G^{\dagger N}\phi \left| \left\{
h_{0}+\sum_{i,j=1}^{r}h_{ij}b_{i}^{\dagger }b_{j}+\frac{1}{4}%
\sum_{i,j,k,l=1}^{r}\left\langle ij||kl\right\rangle b_{i}^{\dagger
}b_{j}^{\dagger }b_{l}b_{k}\right\} b_{p}^{\dagger }b_{p+s}^{\dagger
}G^{\dagger N-1}\phi \right. \right\rangle
\end{eqnarray*}%
\begin{eqnarray}
&=&\sum \zeta _{j_{1}}^{\ast }\cdots \zeta _{j_{N}}^{\ast }\zeta
_{k_{2}}\cdots \zeta _{k_{N}}\qquad \qquad \qquad \qquad \qquad \qquad
\qquad \qquad \qquad \qquad \qquad \quad  \notag \\
&&\qquad \qquad \quad \qquad \langle \eta _{j_{1}}\eta _{j_{1}+s}\cdots \eta
_{j_{N}\,}\eta _{j_{N}+s}\left| H\eta _{p}\eta _{p+s}\eta _{k_{2}}\eta
_{k_{2}+s}\cdots \eta _{k_{N}\,}\eta _{k_{N}+s}\right\rangle  \notag \\
&&\qquad \qquad \qquad \qquad \qquad \qquad \qquad \qquad _{1\leq
j_{1},\cdots ,j_{N}\leq s;\,1\leq k_{2},\cdots ,k_{N}\leq s;\,\,k_{m}\neq
p;\,2\leq m\leq N}
\end{eqnarray}%
The first term produces 
\begin{eqnarray}
&&h_{0}\sum \zeta _{j_{1}}^{\ast }\cdots \zeta _{j_{N}}^{\ast }\zeta
_{k_{2}}\cdots \zeta _{k_{N}}  \notag \\
&&\qquad \qquad \left\langle \eta _{j_{1}}\eta _{j_{1}+s}\cdots \eta
_{j_{N}\,}\eta _{j_{N}+s}\right. \left| \eta _{p}\eta _{p+s}\eta
_{k_{2}}\eta _{k_{2}+s}\cdots \eta _{k_{N}\,}\eta _{k_{N}+s}\right\rangle 
\notag \\
&&\qquad \qquad \qquad \qquad \qquad \qquad \qquad \qquad _{1\leq
j_{1},\cdots ,j_{N}\leq s;\,1\leq k_{2},\cdots ,k_{N}\leq s;\,\,k_{m}\neq
p;\,2\leq m\leq N}  \notag \\
&&\left. {}\right. =h_{0}N!\left( N-1\right) !\zeta _{p}^{\ast
}S_{N-1}\left( \hat{p}\right)
\end{eqnarray}%
The one electron terms are obtained by noting that $b_{i}^{\dagger }b_{j}$
must annihilate the same GSO from the bra and the ket producing 
\begin{eqnarray}
&&\sum \zeta _{j_{1}}^{\ast }\cdots \zeta _{j_{N}}^{\ast }\zeta
_{k_{2}}\cdots \zeta _{k_{N}}  \notag \\
&&\qquad \left\langle \eta _{j_{1}}\eta _{j_{1}+s}\cdots \eta _{j_{N}\,}\eta
_{j_{N}+s}\left| \sum_{i,j=1}^{r}h_{ij}b_{i}^{\dagger }b_{j}\eta _{p}\eta
_{p+s}\eta _{k_{2}}\eta _{k_{2}+s}\cdots \eta _{k_{N}\,}\eta
_{k_{N}+s}\right. \right\rangle  \notag \\
&&\qquad \qquad \qquad \qquad \qquad \qquad \qquad \qquad _{1\leq
j_{1},\cdots ,j_{N}\leq s;\,1\leq k_{2},\cdots ,k_{N}\leq s;\,\,\,k_{m}\neq
p;\quad 2\leq m\leq N}  \notag \\
&=&N!\left( N-1\right) !\zeta _{p}^{\ast }\left\{ \left\{
h_{pp}+h_{p+sp+s}\right\} S_{N-1}\left( \hat{p}\right) +\sum_{i\neq
p}^{s}\left\{ h_{ii}+h_{i+si+s}\right\} n_{j}S_{N-2}\left( \hat{\imath}\hat{p%
}\right) \right\}
\end{eqnarray}%
With the two electron term 
\begin{eqnarray}
&&\sum \zeta _{j_{1}}^{\ast }\cdots \zeta _{j_{N}}^{\ast }\zeta
_{k_{2}}\cdots \zeta _{k_{N}}  \notag \\
&&\qquad \left\langle \eta _{j_{1}}\eta _{j_{1}+s}\cdots \eta _{j_{N}\,}\eta
_{j_{N}+s}\left| \frac{1}{4}\sum_{i,j,k,l=1}^{r}\left\langle
ij||kl\right\rangle b_{i}^{\dagger }b_{j}^{\dagger }b_{l}b_{k}\eta _{p}\eta
_{p+s}\eta _{k_{2}\,}\eta _{k_{2}+s}\cdots \eta _{k_{N}\,}\eta
_{k_{N}+s}\right. \right\rangle  \notag \\
&&\qquad \qquad \qquad \qquad \qquad \qquad \qquad \qquad _{1\leq
j_{1},\cdots ,j_{N}\leq s;\,1\leq k_{2},\cdots ,k_{N}\leq s;\,\,\,k_{m}\neq
p;\quad 2\leq m\leq N}
\end{eqnarray}%
we can use the two electron operator from eq. $\left( \ref{twoel1}\right) $
and $\left( \ref{twoel2}\right) $. From eq. $\left( \ref{twoel1}\right) $ we
obtain 
\begin{eqnarray}
&&4N!\left( N-1\right) !\left\{ \underset{i}{\underset{}{\zeta _{p}^{\ast
}S_{N-1}\left( \hat{p}\right) \left\langle pp+s||pp+s\right\rangle }}+%
\underset{ii}{\sum_{i\neq j\neq p}^{s}\zeta _{p}^{\ast }\zeta _{i}^{\ast
}\zeta _{j}S_{N-2}\left( \hat{p}\hat{\imath}\hat{\jmath}\right) \left\langle
ii+s||jj+s\right\rangle }\right.  \notag \\
&&\left. {}\right. \qquad \qquad \left. \underset{iii}{\sum_{i\neq
p}^{s}\zeta _{p}^{\ast }n_{i}S_{N-2}\left( \hat{p}\hat{\imath}\right)
\left\langle ii+s||ii+s\right\rangle }+\underset{iv}{\sum_{i\neq p}^{s}\zeta
_{i}^{\ast }S_{N-2}\left( \hat{p}\hat{\imath}\right) \left\langle
ii+s||pp+s\right\rangle }\right\}
\end{eqnarray}%
and from eq. $\left( \ref{twoel2}\right) $%
\begin{equation}
4N!\left( N-1\right) !\left\{ \underset{v}{\sum_{i<j\neq p}^{s}\zeta
_{p}^{\ast }n_{i}n_{j}S_{N-3}\left( \hat{p}\hat{\imath}\hat{\jmath}\right)
\left\langle ij|||ij\right\rangle }+\underset{vi}{\sum_{i\neq
p}^{s}n_{i}\zeta _{p}^{\ast }S_{N-2}\left( \hat{p}\hat{\imath}\right)
\left\langle ip|||ip\right\rangle }\right\}
\end{equation}%
Term $\left( i\right) $ comes from summing terms of the form%
\begin{eqnarray}
&&\left\langle pp+s||pp+s\right\rangle \times  \notag \\
&&\qquad \left\langle \eta _{p\,}\eta _{p+s}\eta _{j_{2}}\eta _{j_{2}+s}%
{\small \cdots }\eta _{j_{N}\,}\eta _{j_{N}+s}|b_{p}^{\dagger
}b_{p+s}^{\dagger }b_{p+s}b_{p}\eta _{p}\eta _{p+s}\eta _{j_{2}\,}\eta
_{j_{2}+s}{\small \cdots }\eta _{j_{N}\,}\eta _{j_{N}\,+s}\right\rangle
\label{t1}
\end{eqnarray}%
term $\left( ii\right) $ from%
\begin{eqnarray}
&&\left\langle ii+s||jj+s\right\rangle \times  \notag \\
&&\left\langle \eta _{p\,}\eta _{p+s}\eta _{i}\eta _{i+s}\eta _{j_{3}\,}\eta
_{j_{3}+s}{\small \cdots }\eta _{j_{N}\,}\eta _{j_{N}+s}|b_{i}^{\dagger
}b_{i+s}^{\dagger }b_{j+s}b_{j}\eta _{p}\eta _{p+s}\eta _{j}\eta _{j+s}\eta
_{j_{3}\,}\eta _{j_{3}+s}{\small \cdots }\eta _{j_{N}\,}\eta
_{j_{N}+s}\right\rangle  \label{t2}
\end{eqnarray}%
term $\left( iii\right) $ from%
\begin{eqnarray}
&&\left\langle ii+s||ii+s\right\rangle \times  \notag \\
&&\left\langle \eta _{p\,}\eta _{p+s}\eta _{i}\eta _{i+s}\eta _{j_{3}\,}\eta
_{j_{3}+s}{\small \cdots }\eta _{j_{N}\,}\eta _{j_{N}+s}|b_{i}^{\dagger
}b_{i+s}^{\dagger }b_{i+s}b_{i}\eta _{p}\eta _{p+s}\eta _{i}\eta _{i+s}\eta
_{j_{3}\,}\eta _{j_{3}+s}{\small \cdots }\eta _{j_{N}\,}\eta
_{j_{N}+s}\right\rangle  \label{t3}
\end{eqnarray}%
term $\left( iv\right) $ from%
\begin{equation}
\left\langle ii+s||pp+s\right\rangle \left\langle \eta _{p\,}\eta _{p+s}\eta
_{j_{2}}\eta _{j_{2}+s}\cdots \eta _{j_{N}\,}\eta _{j_{N}+s}|b_{i}^{\dagger
}b_{i+s}^{\dagger }b_{p+s}b_{p}\eta _{p}\eta _{p+s}\eta _{j_{2}\,}\eta
_{j_{2}+s}\cdots \eta _{j_{N}\,}\eta _{j_{N}+s}\right\rangle  \label{t4}
\end{equation}%
term $\left( v\right) $ from%
\begin{eqnarray}
&&\left\langle ij|||ij\right\rangle \times  \notag \\
&&\left\langle \eta _{p\,}\eta _{p+s}\eta _{i}\eta _{i+s}\eta _{j_{3}}\eta
_{j_{3}+s}{\small \cdots }\eta _{j_{N}\,}\eta _{j_{N}+s}|b_{i}^{\dagger
}b_{i+s}^{\dagger }b_{j+s}b_{j}\eta _{p}\eta _{p+s}\eta _{j}\eta _{j+s}\eta
_{j_{3}\,}\eta _{j_{3}+s}{\small \cdots }\eta _{j_{N}\,}\eta
_{j_{N}+s}\right\rangle  \label{t5}
\end{eqnarray}%
and term $\left( vi\right) $ from%
\begin{equation}
\left\langle ip|||ip\right\rangle \left\langle \eta _{p\,}\eta _{p+s}\eta
_{i}\eta _{i+s}\eta _{j_{3}}\eta _{j_{3}+s}\cdots \eta _{j_{N}\,}\eta
_{j_{N}+s}|b_{i}^{\dagger }b_{i+s}^{\dagger }b_{p+s}b_{p}\eta _{p}\eta
_{p+s}\eta _{j_{3}\,}\eta _{j_{3}+s}\cdots \eta _{j_{N}\,}\eta
_{j_{N}+s}\right\rangle  \label{t6}
\end{equation}%
The final result is%
\begin{eqnarray}
&&\left\langle G^{\dagger N}\phi |Hb_{p}^{\dagger }b_{p+s}^{\dagger
}G^{\dagger N-1}\phi \right\rangle =N!\left( N-1\right) !\times  \notag \\
&&\left\{ h_{0}\zeta _{p}^{\ast }S_{N-1}\left( \hat{p}\right) +\right.
\left\{ h_{pp}+h_{p+sp+s}\right\} \zeta _{p}^{\ast }S_{N-1}\left( \hat{p}%
\right) +\sum_{i\neq p}^{s}\left\{ h_{ii}+h_{i+si+s}\right\} \zeta
_{p}^{\ast }n_{i}S_{N-2}\left( \hat{\imath}\hat{p}\right) +  \notag \\
&&\quad \quad \zeta _{p}^{\ast }S_{N-1}\left( \hat{p}\right) \left\langle
pp+s||pp+s\right\rangle +\sum_{i\neq j\neq p}^{s}\zeta _{p}^{\ast }\zeta
_{i}^{\ast }\zeta _{j}S_{N-2}\left( \hat{p}\hat{\imath}\hat{\jmath}\right)
\left\langle ii+s||jj+s\right\rangle  \notag \\
&&\left. {}\right. \qquad \qquad \sum_{i\neq p}^{s}\zeta _{p}^{\ast
}n_{i}S_{N-2}\left( \hat{p}\hat{\imath}\right) \left\langle
ii+s||ii+s\right\rangle +\sum_{i\neq p}^{s}\zeta _{i}^{\ast }S_{N-2}\left( 
\hat{p}\hat{\imath}\right) \left\langle ii+s||pp+s\right\rangle  \notag \\
&&\qquad \qquad \qquad \qquad \left. \sum_{i,j\neq p}^{s}\zeta _{p}^{\ast
}n_{i}n_{j}S_{N-3}\left( \hat{p}\hat{\imath}\hat{\jmath}\right) \left\langle
ij|||ij\right\rangle +\sum_{i\neq p}^{s}n_{i}\zeta _{p}^{\ast }S_{N-2}\left( 
\hat{p}\hat{\imath}\right) \left\langle ip|||ip\right\rangle \right\}
\label{Der}
\end{eqnarray}%
The expression eq. $\left( \ref{Der}\right) $ can be rearranged into sums
driven by the symmetric functions as 
\begin{eqnarray}
&&\left\langle G^{\dagger N}\phi |Hb_{p}^{\dagger }b_{p+s}^{\dagger
}G^{\dagger N-1}\phi \right\rangle =  \notag \\
&&N!\left( N-1\right) !\left\{ \zeta _{p}^{\ast }S_{N-1}\left( \hat{p}%
\right) \left\{ h_{0}+h_{pp}+h_{p+sp+s}+\left\langle pp+s||pp+s\right\rangle
\right\} \right.  \notag \\
&&\left. {}\right. \quad +\sum_{i\neq p}^{s}S_{N-2}\left( \hat{p}\hat{\imath}%
\right) \left\{ \zeta _{p}^{\ast }n_{i}\left\{
h_{ii}+h_{i+si+s}+\left\langle ii+s||ii+s\right\rangle +\left\langle
ip|||ip\right\rangle \right\} \right.  \notag \\
&&\qquad \qquad \qquad \left. +\zeta _{i}^{\ast }\left\langle
ii+s||pp+s\right\rangle \right\} +\sum_{i\neq j\neq p}^{s}\zeta _{i}^{\ast
}\zeta _{j}S_{N-2}\left( \hat{p}\hat{\imath}\hat{\jmath}\right) \left\langle
ii+s||jj+s\right\rangle  \notag \\
&&\qquad \qquad \qquad \qquad \qquad \qquad \qquad \qquad \qquad \left.
+\zeta _{p}^{\ast }\sum_{i<j\neq p}^{s}n_{i}n_{j}S_{N-3}\left( \hat{p}\hat{%
\imath}\hat{\jmath}\right) \left\langle ij|||ij\right\rangle \right\}
\end{eqnarray}

\subsection{\textit{Calculation of} $\left\langle G^{\dagger N}\protect\phi %
|Hb_{p}^{\dagger }b_{q}^{\dagger }G^{\dagger N-1}\protect\phi \right\rangle $%
, $\left| p\right| \neq \left| q\right| $}

In this case there are no contributions from $h_{0}$. The one electron terms
are obtained by noting that $b_{i}^{\dagger }b_{j}$ must either annihilate
the $p$ from the bra and $\bar{q}$ from the ket or $q$ from the bra and $%
\bar{p}$ from the ket to obtain 
\begin{eqnarray}
&&\sum \zeta _{j_{1}}^{\ast }\cdots \zeta _{j_{N}}^{\ast }\zeta
_{k_{2}}\cdots \zeta _{k_{N}}\times  \notag \\
&&\qquad \left\langle \eta _{j_{1}}\eta _{j_{1}+s}\cdots \eta _{j_{N}\,}\eta
_{j_{N}+s}\left| \sum_{i,j=1}^{r}h_{ij}b_{i}^{\dagger }b_{j}\eta _{p}\eta
_{q}\eta _{k_{2}}\eta _{k_{2}+s}\cdots \eta _{k_{N}\,}\eta _{k_{N}+s}\right.
\right\rangle  \notag \\
&&\qquad \qquad \qquad \qquad \qquad \qquad \qquad \qquad _{1\leq
j_{1},\cdots ,j_{N}\leq s;\,1\leq k_{2},\cdots ,k_{N}\leq s;\,\,\,k_{m}\neq
p;\quad 2\leq m\leq N}  \notag \\
&&\left. {}\right. \qquad \qquad \qquad \qquad \qquad \qquad =N!\left(
N-1\right) !\left\{ \zeta _{q}^{\ast }h_{\bar{q}p}+\zeta _{p}^{\ast }h_{\bar{%
p}q}\right\} S_{N-1}\left( \hat{p}\hat{q}\right)
\end{eqnarray}%
The two electron term has the form 
\begin{eqnarray}
&&\sum \zeta _{j_{1}}^{\ast }\cdots \zeta _{j_{N}}^{\ast }\zeta
_{k_{2}}\cdots \zeta _{k_{N}}\times  \notag \\
&&\qquad \left\langle \eta _{j_{1}}\eta _{j_{1}+s}\cdots \eta _{j_{N}\,}\eta
_{j_{N}+s}\left| \frac{1}{4}\sum_{i,j,k,l=1}^{r}\left\langle
ij||kl\right\rangle b_{i}^{\dagger }b_{j}^{\dagger }b_{l}b_{k}\eta _{p}\eta
_{q}\eta _{k_{2}\,}\eta _{k_{2}+s}\cdots \eta _{k_{N}\,}\eta
_{k_{N}+s}\right. \right\rangle  \notag \\
&&\qquad \qquad \qquad \qquad \qquad \qquad \qquad \qquad _{1\leq
j_{1},\cdots ,j_{N}\leq s;\,1\leq k_{2},\cdots ,k_{N}\leq s;\,\,\,k_{m}\neq
p,q;\quad 2\leq m\leq N}
\end{eqnarray}%
which can be expanded as 
\begin{eqnarray*}
&=&\frac{1}{4}\sum_{\kappa ,j_{2},\cdots ,j_{N}\,}\sum_{i,j,k,l=1}^{r}\zeta
_{\kappa }^{\ast }\zeta _{j_{2}}^{\ast }\cdots \zeta _{j_{N}}^{\ast }\zeta
_{j_{2}}\cdots \zeta _{j_{N}}\left\langle \kappa \kappa +s||pq\right\rangle
\times \qquad \qquad \qquad \qquad \qquad \\
&&\left\langle \eta _{\kappa }\eta _{\kappa +s}\eta _{j_{2}}\eta
_{j_{2}+s}\cdots \eta _{j_{N}\,}\eta _{j_{N}+s}\left| b_{\kappa }^{\dagger
}b_{\kappa +s}^{\dagger }b_{q}b_{p}\eta _{p}\eta _{q}\eta _{j_{2}\,}\eta
_{j_{2}+s}\cdots \eta _{j_{N}\,}\eta _{j_{N}+s}\right. \right\rangle
\end{eqnarray*}%
\begin{eqnarray*}
&&+\,\frac{1}{4}\sum_{\kappa ,j_{3},\cdots
,j_{N}\,}\sum_{i,j,k,l=1}^{r}\zeta _{p}^{\ast }\zeta _{q}^{\ast }\zeta
_{j_{3}}^{\ast }\cdots \zeta _{j_{N}}^{\ast }\zeta _{\kappa }\zeta
_{j_{3}}\cdots \zeta _{j_{N}}\left\langle \bar{p}\bar{q}||\kappa \kappa
+s\right\rangle \times \\
&&\left\langle \eta _{p}\eta _{p+s}\eta _{q}\eta _{q+s}\eta _{j_{3}}\eta
_{j_{3}+s}\cdots \eta _{j_{N}\,}\eta _{j_{N}+s}\left| b_{\bar{p}}^{\dagger
}b_{\bar{q}}^{\dagger }b_{\lambda +s}b_{\lambda }\eta _{p}\eta _{q}\eta
_{\lambda \,}\eta _{\lambda +s}\eta _{j_{3}\,}\eta _{j_{3}+s}\cdots \eta
_{j_{N}\,}\eta _{j_{N}+s}\right. \right\rangle
\end{eqnarray*}%
\begin{eqnarray*}
&&+\,\frac{1}{4}\sum_{\kappa ,j_{3},\cdots
,j_{N}\,}\sum_{i,j,k,l=1}^{r}\zeta _{\kappa }^{\ast }\zeta _{q}^{\ast }\zeta
_{j_{3}}^{\ast }\cdots \zeta _{j_{N}}^{\ast }\zeta _{j_{3}}\cdots \zeta
_{j_{N}}\left\langle \bar{q}\kappa ||p\kappa \right\rangle \times \\
&&\left\langle \eta _{\kappa }\eta _{\kappa +s}\eta _{q}\eta _{q+s}\eta
_{j_{3}\,}\eta _{j_{3}+s}\cdots \eta _{j_{N}\,}\eta _{j_{N}+s}\left| b_{\bar{%
q}}^{\dagger }b_{\kappa }^{\dagger }b_{\kappa }b_{p}\eta _{p}\eta _{q}\eta
_{\kappa \,}\eta _{\kappa +s}\eta _{j_{3}\,}\eta _{j_{3}+s}\cdots \eta
_{j_{N}\,}\eta _{j_{N}+s}\right. \right\rangle
\end{eqnarray*}%
\begin{eqnarray}
&&+\,\frac{1}{4}\sum_{\kappa ,j_{3},\cdots
,j_{N}\,}\sum_{i,j,k,l=1}^{r}\zeta _{\kappa }^{\ast }\zeta _{q}^{\ast }\zeta
_{j_{3}}^{\ast }\cdots \zeta _{j_{N}}^{\ast }\zeta _{j_{3}}\cdots \zeta
_{j_{N}}\left\langle \bar{q}\kappa ||p\kappa \right\rangle \times  \notag \\
&&\left\langle \eta _{\kappa }\eta _{\kappa +s}\eta _{q}\eta _{q+s}\eta
_{j_{3}\,}\eta _{j_{3}+s}\cdots \eta _{j_{N}\,}\eta _{j_{N}+s}\left| b_{\bar{%
q}}^{\dagger }b_{\kappa }^{\dagger }b_{\kappa }b_{p}\eta _{p}\eta _{q}\eta
_{\kappa \,}\eta _{\kappa +s}\eta _{j_{3}\,}\eta _{j_{3}+s}\cdots \eta
_{j_{N}\,}\eta _{j_{N}+s}\right. \right\rangle
\end{eqnarray}%
which gives 
\begin{eqnarray}
&&\sum_{\kappa \neq \left| p\right| \neq \left| q\right| }^{s}\left\{ \zeta
_{\kappa }^{\ast }\left\langle \kappa \kappa +s||pq\right\rangle
S_{N-1}\left( \hat{\kappa}\hat{p}\hat{q}\right) +\right.  \notag \\
&&\qquad \qquad \left. +\zeta _{p}^{\ast }\left\langle pp+s||pq\right\rangle
S_{N-1}\left( \hat{p}\hat{q}\right) +\zeta _{q}^{\ast }\left\langle
qq+s||pq\right\rangle S_{N-1}\left( \hat{p}\hat{q}\right) \right\}  \notag \\
&&\qquad \qquad +\sum_{\kappa \neq \left| p\right| \neq \left| q\right|
}^{s}\left\langle \bar{p}\bar{q}||\kappa \kappa +s\right\rangle \zeta
_{p}^{\ast }\zeta _{q}^{\ast }\zeta _{\kappa }S_{N-2}\left( \hat{\kappa}\hat{%
p}\hat{q}\right)  \notag \\
&&\qquad \qquad +\sum_{\kappa \neq \left| p\right| \neq \left| q\right|
}^{s}\left\langle \bar{q}\kappa ||p\kappa \right\rangle n_{\kappa }\zeta
_{q}^{\ast }S_{N-2}\left( \hat{\kappa}\hat{p}\hat{q}\right)  \notag \\
&&\qquad \qquad +\sum_{\kappa \neq \left| p\right| \neq \left| q\right|
}^{s}\left\langle \bar{p}\kappa ||q\kappa \right\rangle n_{\kappa }\zeta
_{p}^{\ast }S_{N-2}\left( \hat{\kappa}\hat{p}\hat{q}\right)
\end{eqnarray}%
which can be collected together with the one electron terms to give 
\begin{eqnarray}
&&\left\langle G^{\dagger N}\phi |Hb_{p}^{\dagger }b_{q}^{\dagger
}G^{\dagger N-1}\phi \right\rangle =N!\left( N-1\right) !\times  \notag \\
&&\left\{ \sum_{\kappa \neq \left| p\right| \neq \left| q\right|
}^{s}\left\{ S_{N-1}\left( \hat{\kappa}\hat{p}\hat{q}\right) \zeta _{\kappa
}^{\ast }\left\langle \kappa \kappa +s||pq\right\rangle +S_{N-2}\left( \hat{%
\kappa}\hat{p}\hat{q}\right) \left\langle \bar{p}\bar{q}||\kappa \kappa
+s\right\rangle \zeta _{p}^{\ast }\zeta _{q}^{\ast }\zeta _{\kappa }\right.
\right.  \notag \\
&&\quad \quad \quad \quad \quad \quad \left. +\left\langle \bar{q}\kappa
||p\kappa \right\rangle n_{\kappa }\zeta _{q}^{\ast }+\left\langle \bar{p}%
\kappa ||q\kappa \right\rangle n_{\kappa }\zeta _{p}^{\ast }\right\} +\zeta
_{p}^{\ast }\left\langle pp+s||pq\right\rangle S_{N-1}\left( \hat{p}\hat{q}%
\right)  \notag \\
&&\quad \quad \quad \quad \quad \quad \quad \left. +\zeta _{q}^{\ast
}\left\langle qq+s||pq\right\rangle S_{N-1}\left( \hat{p}\hat{q}\right)
+\left\{ \zeta _{q}^{\ast }h_{\bar{q}p}+\zeta _{p}^{\ast }h_{\bar{p}%
q}\right\} S_{N-1}\left( \hat{p}\hat{q}\right) \overset{}{\underset{}{}}%
\right\}
\end{eqnarray}

\section{Transformation of the Energy of the 2N-Electron Unnormalized AGP
State to the Atomic Orbital Basis\label{Trans2AObasis}}

Expressing the canonical basis $\left\{ \eta _{j}\right\} $ in terms of the
atomic basis $\left\{ \varphi _{j}\right\} $ the energy expression eq. $%
\left( \ref{canen}\right) $ becomes

\begin{eqnarray}
&&\mathcal{H}\left( \mathbf{\zeta }^{\ast },\mathbf{\zeta }\right) =h_{0}%
\mathcal{N}\left( g^{N}\right) +\left( N\right) !^{2}\times  \notag \\
&&\left\{ \sum_{j_{1,}j_{2}=1}^{r}\sum_{1\leq j\leq s}n_{j}S_{N-1}\left( 
\hat{\jmath}\right) \left\{ \underset{}{\overset{}{}}\left\{
V_{f\,j_{1}j}^{\ast }V_{f\,j_{2}j}+V_{f\,j_{1}j+s}^{\ast
}V_{f\,j_{2}j+s}\right\} h_{j_{1}j_{2}}^{a}\right. \right.  \notag \\
&&+\sum_{j_{1,}j_{2,}j_{3,}j_{4}=1}^{r}V_{f\,j_{1}j}^{\ast
}V_{f\,j_{2}j+s}^{\ast }V_{f\,j_{3}j}V_{f\,j_{4}j+s}\left\langle
j_{1}j_{2}||j_{3}j_{4}\right\rangle ^{a}\left. \underset{}{\overset{}{}}%
\right\}  \notag \\
&&+\sum_{j_{1,}j_{2,}j_{3,}j_{4}=1}^{r}\sum_{1\leq j\neq k\leq s}\zeta
_{j}^{\ast }\zeta _{k}S_{N-1}\left( \hat{\jmath}\hat{k}\right)
V_{f\,j_{1}j}^{\ast }V_{f\,j_{2}j+s}^{\ast
}V_{f\,j_{3}k}V_{f\,j_{4}k+s}\left\langle
j_{1}j_{2}||j_{3}j_{4}\right\rangle ^{a}  \notag \\
&&\left. +\sum_{j_{1,}j_{2,}j_{3,}j_{4}=1}^{r}\sum_{1\leq j,k\leq
s}n_{j}n_{k}S_{N-2}\left( \hat{\jmath}\hat{k}\right) V_{f\,j_{1}j}^{\ast
}V_{f\,j_{2}k}^{\ast }V_{f\,j_{3}j}V_{f\,j_{4}k}\left\langle
j_{1}j_{2}||j_{3}j_{4}\right\rangle ^{a}\right\}
\end{eqnarray}%
which can be rearranged to give 
\begin{eqnarray}
&&\mathcal{H}\left( \mathbf{\zeta }^{\ast },\mathbf{\zeta }\right) =h_{0}%
\mathcal{N}\left( g^{N}\right) +\left( N\right) !^{2}\times  \notag \\
&&\left\{ \sum_{j_{1,}j_{2}=1}^{r}\sum_{1\leq j\leq s}n_{j}S_{N-1}\left( 
\hat{\jmath}\right) \left\{ V_{f\,j_{1}j}^{\ast
}V_{f\,j_{2}j}+V_{f\,j_{1}j+s}^{\ast }V_{f\,j_{2}j+s}\right\}
h_{j_{1}j_{2}}^{a}\right.  \notag \\
&&+\sum_{j_{1,}j_{2,}j_{3,}j_{4}=1}^{r}\left\{ \overset{}{\underset{}{}}%
\sum_{1\leq j\leq s}n_{j}S_{N-1}\left( \hat{\jmath}\right)
V_{f\,j_{1}j}^{\ast }V_{f\,j_{2}j+s}^{\ast
}V_{f\,j_{3}j}V_{f\,j_{4}j+s}\right.  \notag \\
&&+\sum_{1\leq j\neq k\leq s}\zeta _{j}^{\ast }\zeta _{k}S_{N-1}\left( \hat{%
\jmath}\hat{k}\right) V_{f\,j_{1}j}^{\ast }V_{f\,j_{2}j+s}^{\ast
}V_{f\,j_{3}k}V_{f\,j_{4}k+s}  \notag \\
&&\left. \left. +\sum_{1\leq j,k\leq s}n_{j}n_{k}S_{N-2}\left( \hat{\jmath}%
\hat{k}\right) V_{f\,j_{1}j}^{\ast }V_{f\,j_{2}k}^{\ast
}V_{f\,j_{3}j}V_{f\,j_{4}k}\right\} \left\langle
j_{1}j_{2}||j_{3}j_{4}\right\rangle ^{a}\right\}
\end{eqnarray}

\section{Spin Integration of the Hamiltonian Matrix Elements \label%
{SpinInteg}}

\begin{eqnarray}
&&h_{jk}^{a}=h_{j+sk+s}^{a}=-\int \chi _{j}^{\ast }\left( \mathbf{r}\right)
\left\{ \sum_{_{1\leq j\leq N_{nuc}}}\frac{Z_{j}}{\left\| \mathbf{R}_{j}-%
\mathbf{r}\right\| }+\frac{1}{2m_{e}}\nabla _{\mathbf{r}}^{2}\right\} \chi
_{k}\left( \mathbf{r}\right) d\mathbf{r};\quad 1\leq j,k\leq s  \notag \\
&&h_{jk+s}^{a}=h_{j+sk}^{a}=0;\quad 1\leq j,k\leq s \\
&&  \notag \\
&&\left\langle jk||lm\right\rangle ^{a}=\int \left\{ \chi _{j}^{\ast }\left( 
\mathbf{r}_{1}\right) \chi _{k}^{\ast }\left( \mathbf{r}_{2}\right) \frac{1}{%
\left\| \mathbf{r}_{1}-\mathbf{r}_{2}\right\| }\chi _{l}\left( \mathbf{r}%
_{1}\right) \chi _{m}\left( \mathbf{r}_{2}\right) \right.  \notag \\
&&-\left. \chi _{j}^{\ast }\left( \mathbf{r}_{1}\right) \chi _{k}^{\ast
}\left( \mathbf{r}_{2}\right) \frac{1}{\left\| \mathbf{r}_{1}-\mathbf{r}%
_{2}\right\| }\chi _{l}\left( \mathbf{r}_{2}\right) \chi _{m}\left( \mathbf{r%
}_{1}\right) \right\} d\mathbf{r}_{1}d\mathbf{r}_{2};\quad 1\leq j,k\leq
s;1\leq l,m\leq s \\
&&  \notag \\
&&\left\langle j+sk||lm\right\rangle ^{a}=\left\langle jk+s||lm\right\rangle
^{a}=\left\langle jk||l+sm\right\rangle ^{a}=\left\langle
jk||lm+s\right\rangle ^{a}=\left\langle j+sk+s||l+sm\right\rangle ^{a}= 
\notag \\
&&\left\langle j+sk+s||lm+s\right\rangle ^{a}=\left\langle
j+sk||l+sm+s\right\rangle ^{a}=\left\langle jk+s||l+sm+s\right\rangle ^{a}=0
\notag \\
&&\left. {}\right. \quad \qquad \qquad \quad \qquad \qquad \quad \qquad
\qquad \quad \qquad \qquad \quad 1\leq j,k\leq s;1\leq l,m\leq s \\
&&  \notag \\
&&\left\langle jk+s||l+sm\right\rangle ^{a}=\left\langle
j+sk||l+sm\right\rangle ^{a}=  \notag \\
&&\qquad \int \chi _{j}^{\ast }\left( \mathbf{r}_{1}\right) \chi _{k}^{\ast
}\left( \mathbf{r}_{2}\right) \frac{1}{\left\| \mathbf{r}_{1}-\mathbf{r}%
_{2}\right\| }\chi _{l}\left( \mathbf{r}_{1}\right) \chi _{m}\left( \mathbf{r%
}_{2}\right) d\mathbf{r}_{1}d\mathbf{r}_{2};\quad 1\leq j,k\leq s;1\leq
l,m\leq s \\
&&  \notag \\
&&\left\langle jk+s||lm+s\right\rangle ^{a}=\left\langle
j+sk||l+sm\right\rangle ^{a}=-\left\langle jk+s||l+sm\right\rangle
^{a}=-\left\langle j+sk||lm+s\right\rangle ^{a}  \notag \\
&&\left. {}\right. \quad \qquad \qquad \quad \qquad \qquad \quad \quad
\qquad \quad \qquad \qquad \qquad \quad 1\leq j,k\leq s;1\leq l,m\leq s \\
&&\left\langle jk|||jk\right\rangle ^{a}=\left\langle jk||jk\right\rangle
^{a}+\left\langle j\,k+s||j\,k+s\right\rangle ^{a}+\left\langle
j+s\,k||j+s\,k\right\rangle ^{a}+\left\langle
j+s\,k+s||j+s\,k+s\right\rangle ^{a}  \notag \\
&&\left. {}\right. \quad \qquad \qquad \quad \qquad \qquad \quad \qquad
\qquad \quad \qquad =0;\qquad 1\leq j,k\leq s;1\leq l,m\leq s
\end{eqnarray}

\section{Transformation of Derivatives to Atomic Orbital Basis \label%
{DerTrans}}

In the atomic orbital case we only have derivatives of the form%
\begin{equation}
\left\langle G^{\dagger N}\phi |Ha_{j_{1}}^{\dagger }a_{j_{2}}^{\dagger
}G^{\dagger N-1}\phi \right\rangle ;1\leq j_{1},j_{2}\leq r
\end{equation}%
which are given in terms of the derivatives expressed in the canonical basis
as%
\begin{equation}
\left\langle G^{\dagger N}\phi |Ha_{j_{1}}^{\dagger }a_{j_{2}}^{\dagger
}G^{\dagger N-1}\phi \right\rangle
=\sum_{p,q=1}^{r}V_{f\,j_{1}p}V_{f\,j_{2}q}\left\langle G^{\dagger N}\phi
|Hb_{p}^{\dagger }b_{q}^{\dagger }G^{\dagger N-1}\phi \right\rangle
\end{equation}%
a sum that goes over the two types of matrix elements computed in eq. $%
\left( \ref{pps}\right) $ and eq. $\left( \ref{pq}\right) $ and each of
these matrix elements can also be expanded in the atomic basis. The paired
terms $p=\left| q\right| $ as%
\begin{eqnarray*}
&&\left\langle G^{\dagger N}\phi |Hb_{p}^{\dagger }b_{p+s}^{\dagger
}G^{\dagger N-1}\phi \right\rangle =N!\left( N-1\right) !\times \qquad
\qquad \qquad \qquad \qquad \qquad \\
&&\left\{ S_{N-1}\left( \hat{p}\right) \zeta _{p}^{\ast }\left\{
h_{0}+\sum_{j_{1,}j_{2}=1}^{r}\left\{ V_{f\,j_{1}p}^{\ast
}V_{f\,j_{2}p}+V_{f\,j_{1}p+s}^{\ast }V_{f\,j_{2}p+s}\right\}
h_{j_{1}j_{2}}^{a}\right. \right. \\
&&+\left. {}\right. \qquad \qquad \quad \left.
\sum_{j_{1,}j_{2,}j_{3,}j_{4}=1}^{r}V_{f\,j_{1}p}^{\ast
}V_{f\,j_{2}p+s}^{\ast }V_{f\,j_{3}p}V_{f\,j_{4}p+s}\left\langle
j_{1}j_{2}||j_{3}j_{4}\right\rangle ^{a}\right\}
\end{eqnarray*}%
\begin{eqnarray*}
&&+\sum_{j\neq p}^{s}S_{N-2}\left( \hat{p}\hat{\jmath}\right) \left\{ \zeta
_{p}^{\ast }n_{j}\left\{ \sum_{j_{1},j_{2}}^{r}\left\{ V_{f\,j_{1}j}^{\ast
}V_{f\,j_{2}j}+V_{f\,j_{1}j+s}^{\ast }V_{f\,j_{2}j+s}\right\}
h_{j_{1}j_{2}}^{a}\right. \right. \qquad \\
&&+\left. \sum_{j_{1,}j_{2,}j_{3,}j_{4}=1}^{r}\left\{ V_{f\,j_{1}j}^{\ast
}V_{f\,j_{2}j+s}^{\ast }V_{f\,j_{3}j}V_{f\,j_{4}j+s}+V_{f\,j_{1}j}^{\ast
}V_{f\,j_{2}p}^{\ast }V_{f\,j_{3}p}V_{f\,j_{4}j}\right\} \left\langle
j_{1}j_{2}||j_{3}j_{4}\right\rangle ^{a}\right\}
\end{eqnarray*}%
\begin{eqnarray*}
&&+\sum_{j_{1,}j_{2,}j_{3,}j_{4}=1}^{r}\left\{ V_{f\,j_{1}j}^{\ast
}V_{f\,j_{2}p}^{\ast }V_{f\,j_{3}p}V_{f\,j_{4}j}+V_{f\,j_{1}j}^{\ast
}V_{f\,j_{2}p+s}^{\ast }V_{f\,j_{3}p+s}V_{f\,j_{4}j}\right. \\
&&\qquad \qquad \qquad \qquad +V_{f\,j_{1}j+s}^{\ast }V_{f\,j_{2}p}^{\ast
}V_{f\,j_{3}p}V_{f\,j_{4}j+s}+V_{f\,j_{1}j+s}^{\ast }V_{f\,j_{2}p+s}^{\ast
}V_{f\,j_{3}p+s}V_{f\,j_{4}j+s} \\
&&\left. \qquad \qquad \qquad \qquad \qquad \qquad \qquad \right. +\left.
\zeta _{j}^{\ast }V_{f\,j_{1}j}^{\ast }V_{f\,j_{2}j+s}^{\ast
}V_{f\,j_{3}p}V_{f\,j_{4}p+s}\right\} \left\langle
j_{1}j_{2}||j_{3}j_{4}\right\rangle ^{a}
\end{eqnarray*}%
\begin{eqnarray}
&&+\sum_{j\neq k\neq p}S_{N-2}\left( \hat{p}\hat{\jmath}\hat{k}\right) \zeta
_{p}^{\ast }\zeta _{j}^{\ast }\zeta
_{k}\sum_{j_{1,}j_{2,}j_{3,}j_{4}=1}^{r}V_{f\,j_{1}j}^{\ast
}V_{f\,j_{2}j+s}^{\ast }V_{f\,j_{3}k}V_{f\,j_{4}k+s}\left\langle
j_{1}j_{2}||j_{3}j_{4}\right\rangle ^{a}  \notag \\
&&+\frac{1}{2}\zeta _{p}^{\ast }\sum_{j,k\neq p}S_{N-3}\left( \hat{p}\hat{%
\jmath}\hat{k}\right) n_{j}n_{k}\sum_{j_{1,}j_{2,}j_{3,}j_{4}=1}^{r}\left\{
V_{f\,j_{1}j}^{\ast }V_{f\,j_{2}k}^{\ast }V_{f\,j_{3}k}V_{f\,j_{4}pj}\right.
\notag \\
&&\qquad \qquad \qquad +V_{f\,j_{1}j}^{\ast }V_{f\,j_{2}k+s}^{\ast
}V_{f\,j_{3}k+s}V_{f\,j_{4}pj}+V_{f\,j_{1}j+s}^{\ast }V_{f\,j_{2}k}^{\ast
}V_{f\,j_{3}k}V_{f\,j_{4}pj+s}+  \notag \\
&&\qquad \qquad \qquad \qquad \qquad \quad +\left. \left.
+V_{f\,j_{1}j+s}^{\ast }V_{f\,j_{2}k+s}^{\ast
}V_{f\,j_{3}k+s}V_{f\,j_{4}pj+s}\right\} \left\langle
j_{1}j_{2}||j_{3}j_{4}\right\rangle ^{a}\underset{}{\overset{}{}}\right\}
\end{eqnarray}%
which can be rearranged as%
\begin{eqnarray*}
&&\left\langle G^{\dagger N}\phi |Hb_{p}^{\dagger }b_{p+s}^{\dagger
}G^{\dagger N-1}\phi \right\rangle =N!\left( N-1\right) !\times \\
&&\left\{ h_{0}S_{N-1}\left( \hat{p}\right) \zeta _{p}^{\ast
}+\sum_{j_{1,}j_{2}=1}^{r}\left\{ \overset{}{\underset{}{}}S_{N-1}\left( 
\hat{p}\right) \zeta _{p}^{\ast }\left\{ V_{f\,j_{1}p}^{\ast
}V_{f\,j_{2}p}+V_{f\,j_{1}p+s}^{\ast }V_{f\,j_{2}p+s}\right\}
h_{j_{1}j_{2}}^{a}\right. \right. \\
&&+\left. \sum_{j\neq p}^{s}S_{N-2}\left( \hat{p}\hat{\jmath}\right) \zeta
_{p}^{\ast }n_{j}\left\{ V_{f\,j_{1}j}^{\ast
}V_{f\,j_{2}j}+V_{f\,j_{1}j+s}^{\ast }V_{f\,j_{2}j+s}\right\}
h_{j_{1}j_{2}}^{a}\right\} \qquad \qquad \qquad \qquad \qquad
\end{eqnarray*}%
\begin{eqnarray*}
&&+\sum_{j_{1,}j_{2,}j_{3,}j_{4}=1}^{r}\left\{ \overset{}{\underset{}{}}%
S_{N-1}\left( \hat{p}\right) \zeta _{p}^{\ast }V_{f\,j_{1}p}^{\ast
}V_{f\,j_{2}p+s}^{\ast }V_{f\,j_{3}p}V_{f\,j_{4}p+s}\right. \\
&&+\sum_{j\neq p}^{s}S_{N-2}\left( \hat{p}\hat{\jmath}\right) \left\{ \zeta
_{p}^{\ast }n_{j}\left\{ V_{f\,j_{1}j}^{\ast }V_{f\,j_{2}j+s}^{\ast
}V_{f\,j_{3}j}V_{f\,j_{4}j+s}+V_{f\,j_{1}j}^{\ast }V_{f\,j_{2}p}^{\ast
}V_{f\,j_{3}p}V_{f\,j_{4}j}\right. \right. \\
&&\qquad \qquad \qquad \qquad +V_{f\,j_{1}j}^{\ast }V_{f\,j_{2}p+s}^{\ast
}V_{f\,j_{3}p+s}V_{f\,j_{4}j}+V_{f\,j_{1}j+s}^{\ast }V_{f\,j_{2}p}^{\ast
}V_{f\,j_{3}p}V_{f\,j_{4}j+s} \\
&&\qquad \qquad \qquad \qquad \qquad +\left. V_{f\,j_{1}j+s}^{\ast
}V_{f\,j_{2}p+s}^{\ast }V_{f\,j_{3}p+s}V_{f\,j_{4}j+s}\right\} \left. +\zeta
_{j}^{\ast }V_{f\,j_{1}j}^{\ast }V_{f\,j_{2}j+s}^{\ast
}V_{f\,j_{3}p}V_{f\,j_{4}p+s}\right\}
\end{eqnarray*}%
\begin{eqnarray}
&&+\sum_{j\neq k\neq p}S_{N-2}\left( \hat{p}\hat{\jmath}\hat{k}\right) \zeta
_{p}^{\ast }\zeta _{j}^{\ast }\zeta _{k}V_{f\,j_{1}j}^{\ast
}V_{f\,j_{2}j+s}^{\ast }V_{f\,j_{3}k}V_{f\,j_{4}k+s}\qquad \qquad \qquad 
\notag \\
&&+\frac{1}{2}\zeta _{p}^{\ast }\sum_{j,k\neq p}S_{N-3}\left( \hat{p}\hat{%
\jmath}\hat{k}\right) n_{j}n_{k}\left\{ V_{f\,j_{1}j}^{\ast
}V_{f\,j_{2}k}^{\ast }V_{f\,j_{3}k}V_{f\,j_{4}pj}\right.  \notag \\
&&\qquad \qquad \qquad +V_{f\,j_{1}j}^{\ast }V_{f\,j_{2}k+s}^{\ast
}V_{f\,j_{3}k+s}V_{f\,j_{4}pj}+V_{f\,j_{1}j+s}^{\ast }V_{f\,j_{2}k}^{\ast
}V_{f\,j_{3}k}V_{f\,j_{4}pj+s}  \notag \\
&&\qquad \qquad \qquad \qquad \qquad \qquad \left. \left. \left.
+V_{f\,j_{1}j+s}^{\ast }V_{f\,j_{2}k+s}^{\ast
}V_{f\,j_{3}k+s}V_{f\,j_{4}pj+s}\right\} \underset{}{\overset{}{}}\right\}
\left\langle j_{1}j_{2}||j_{3}j_{4}\right\rangle ^{a}\right\}
\end{eqnarray}%
The other term $\left| p\right| \neq \left| q\right| $%
\begin{eqnarray}
&&\left\langle G^{\dagger N}\phi |Hb_{p}^{\dagger }b_{q}^{\dagger
}G^{\dagger N-1}\phi \right\rangle =N!\left( N-1\right) !\times  \notag \\
&&\left\{ \sum_{\kappa \neq \left| p\right| \neq \left| q\right|
}^{s}\left\{ S_{N-1}\left( \hat{\kappa}\hat{p}\hat{q}\right) \left\{
\sum_{j_{1,}j_{2,}j_{3,}j_{4}=1}^{r}V_{f\,j_{1}\kappa }^{\ast
}V_{f\,j_{2}\kappa +s}^{\ast }V_{f\,j_{3}p}V_{f\,j_{4}q}\zeta _{\kappa
}^{\ast }\left\langle j_{1}j_{2}||j_{3}j_{4}\right\rangle ^{a}\right.
\right. \right.  \notag \\
&&\quad \qquad \quad \quad \quad \qquad +\left.
\sum_{j_{1,}j_{2,}j_{3,}j_{4}=1}^{r}V_{f\,j_{1}\bar{p}}^{\ast }V_{f\,j_{2}%
\bar{q}}^{\ast }V_{f\,j_{3}\kappa }V_{f\,j_{4}\kappa +s}\left\langle
j_{1}j_{2}||j_{3}j_{4}\right\rangle ^{a}\zeta _{p}^{\ast }\zeta _{q}^{\ast
}\zeta _{\kappa }\right\}  \notag \\
&&\qquad \qquad \quad \quad \qquad
+\sum_{j_{1,}j_{2,}j_{3,}j_{4}=1}^{r}V_{f\,j_{1}\bar{q}}^{\ast
}V_{f\,j_{2}\kappa }^{\ast }V_{f\,j_{3}p}V_{f\,j_{4}\kappa }\left\langle
j_{1}j_{2}||j_{3}j_{4}\right\rangle ^{a}n_{\kappa }\zeta _{q}^{\ast }  \notag
\\
&&\qquad \qquad \quad \quad \qquad \left.
+\sum_{j_{1,}j_{2,}j_{3,}j_{4}=1}^{r}V_{f\,j_{1}\bar{p}}^{\ast
}V_{f\,j_{2}\kappa }^{\ast }V_{f\,j_{3}q}V_{f\,j_{4}\kappa }\left\langle
j_{1}j_{2}||j_{3}j_{4}\right\rangle ^{a}n_{\kappa }\zeta _{p}^{\ast }\right\}
\notag \\
&&\qquad \qquad \quad \quad \quad +\zeta _{p}^{\ast }S_{N-1}\left( \hat{p}%
\hat{q}\right) \sum_{j_{1,}j_{2,}j_{3,}j_{4}=1}^{r}V_{f\,j_{1}p}^{\ast
}V_{f\,j_{2}p+s}^{\ast }V_{f\,j_{3}p}V_{f\,j_{4}q}\left\langle
j_{1}j_{2}||j_{3}j_{4}\right\rangle ^{a}  \notag \\
&&\qquad \qquad \quad \quad \quad +\zeta _{q}^{\ast }S_{N-1}\left( \hat{p}%
\hat{q}\right) \sum_{j_{1,}j_{2,}j_{3,}j_{4}=1}^{r}V_{f\,j_{1}q}^{\ast
}V_{f\,j_{2}q+s}^{\ast }V_{f\,j_{3}p}V_{f\,j_{4}q}\left\langle
j_{1}j_{2}||j_{3}j_{4}\right\rangle ^{a}  \notag \\
&&\qquad \quad \quad \quad \qquad \left. +S_{N-1}\left( \hat{p}\hat{q}%
\right) \sum_{j_{1},j_{2}=1}^{r}\left\{ V_{f\,j_{1}\bar{q}}^{\ast
}V_{f\,j_{2}p}\zeta _{q}^{\ast }+V_{f\,j_{1}\bar{p}}^{\ast
}V_{f\,j_{2}q}\zeta _{p}^{\ast }\right\} h_{j_{1}j_{2}}^{a}\right\}
\end{eqnarray}%
This can be rearranged to give 
\begin{eqnarray}
&&\left\langle G^{\dagger N}\phi |Hb_{p}^{\dagger }b_{q}^{\dagger
}G^{\dagger N-1}\phi \right\rangle =N!\left( N-1\right) !\times  \notag \\
&&\left\{ S_{N-1}\left( \hat{p}\hat{q}\right)
\sum_{j_{1},j_{2}=1}^{r}\left\{ V_{f\,j_{1}\bar{q}}^{\ast
}V_{f\,j_{2}p}\zeta _{q}^{\ast }+V_{f\,j_{1}\bar{p}}^{\ast
}V_{f\,j_{2}q}\zeta _{p}^{\ast }\right\} h_{j_{1}j_{2}}\right.  \notag \\
&&+\sum_{j_{1,}j_{2,}j_{3,}j_{4}=1}^{r}\left\{ \sum_{\kappa \neq \left|
p\right| \neq \left| q\right| }^{s}S_{N-1}\left( \hat{\kappa}\hat{p}\hat{q}%
\right) \left\{ V_{f\,j_{1}\kappa }^{\ast }V_{f\,j_{2}\kappa +s}^{\ast
}V_{f\,j_{3}p}V_{f\,j_{4}q}\zeta _{\kappa }^{\ast }\right. \right.  \notag \\
&&\quad \qquad \quad +V_{f\,j_{1}\bar{p}}^{\ast }V_{f\,j_{2}\bar{q}}^{\ast
}V_{f\,j_{3}\kappa }V_{f\,j_{4}\kappa +s}\zeta _{p}^{\ast }\zeta _{q}^{\ast
}\zeta _{\kappa }+V_{f\,j_{1}\bar{q}}^{\ast }V_{f\,j_{2}\kappa }^{\ast
}V_{f\,j_{3}p}V_{f\,j_{4}\kappa }n_{\kappa }\zeta _{q}^{\ast }  \notag \\
&&\qquad \qquad \qquad \qquad \qquad \qquad \qquad \qquad \qquad \left.
+V_{f\,j_{1}\bar{p}}^{\ast }V_{f\,j_{2}\kappa }^{\ast
}V_{f\,j_{3}q}V_{f\,j_{4}\kappa }n_{\kappa }\zeta _{p}^{\ast }\right\} 
\notag \\
&&+\zeta _{p}^{\ast }S_{N-1}\left( \hat{p}\hat{q}\right) V_{f\,j_{1}p}^{\ast
}V_{f\,j_{2}p+s}^{\ast }V_{f\,j_{3}p}V_{f\,j_{4}q}  \notag \\
&&\qquad \qquad \qquad \left. \left. +\zeta _{q}^{\ast }S_{N-1}\left( \hat{p}%
\hat{q}\right) V_{f\,j_{1}q}^{\ast }V_{f\,j_{2}q+s}^{\ast
}V_{f\,j_{3}p}V_{f\,j_{4}q}\underset{}{\overset{}{}}\right\} \left\langle
j_{1}j_{2}||j_{3}j_{4}\right\rangle ^{a}\underset{}{\overset{\underset{}{}}{}%
}\right\}
\end{eqnarray}

\section{Cost of the Computations\label{CostofComp}}

The partial derivatives of the energy wrt geminal coefficients are given by
eq. $\left( \ref{deriv}\right) $ as a sum that goes over two types of matrix
elements, paired terms where $p=\left| q\right| $ and unpaired terms where $%
p\neq \left| q\right| $%
\begin{eqnarray}
&&\left\langle G^{\dagger N}\phi |Hb_{p}^{\dagger }b_{p+s}^{\dagger
}G^{\dagger N-1}\phi \right\rangle =N!\left( N-1\right) !\times  \notag \\
&&\left\{ h_{0}S_{N-1}\left( \hat{p}\right) \zeta _{p}^{\ast
}+\sum_{j_{1,}j_{2}=1}^{r}\left\{ \underset{}{\overset{}{}}S_{N-1}\left( 
\hat{p}\right) \zeta _{p}^{\ast }\left\{ V_{f\,j_{1}p}^{\ast
}V_{f\,j_{2}p}+V_{f\,j_{1}p+s}^{\ast }V_{f\,j_{2}p+s}\right\}
h_{j_{1}j_{2}}^{a}\right. \right.  \notag \\
&&\qquad \qquad \qquad \left. +\right. \left. \sum_{j\neq
p}^{s}S_{N-2}\left( \hat{p}\hat{\jmath}\right) \zeta _{p}^{\ast
}n_{j}\left\{ V_{f\,j_{1}j}^{\ast }V_{f\,j_{2}j}+V_{f\,j_{1}j+s}^{\ast
}V_{f\,j_{2}j+s}\right\} h_{j_{1}j_{2}}^{a}\right\} .
\end{eqnarray}%
The paired terms themselves can be expressed as a sum of three types $%
TP_{0},TP_{1}$ and $TP_{2}$ 
\begin{eqnarray*}
&&\left\langle G^{\dagger N}\phi |Hb_{p}^{\dagger }b_{p+s}^{\dagger
}G^{\dagger N-1}\phi \right\rangle =TP_{0}+TP_{1}+TP_{2} \\
&=&N!\left( N-1\right) !\left\{
TP_{0p}h_{0}+\sum_{j_{1,}j_{2}=1}TP_{1pj_{1}j_{2}}h_{j_{1}j_{2}}^{a}+%
\sum_{j_{1,}j_{2,}j_{3,}j_{4}=1}^{r}TP_{2pj_{1}j_{2}j_{3}j_{4}}\left\langle
j_{1}j_{2}||j_{3}j_{4}\right\rangle ^{a}\right\}
\end{eqnarray*}%
where $TP_{0}$ is given by 
\begin{equation}
TP_{0p}=S_{N-1}\left( \hat{p}\right) \zeta _{p}^{\ast },
\end{equation}%
$TP_{1}$ is a term whose elements' computational cost scale as $s$ and is
given by 
\begin{eqnarray}
&&TP_{1pj_{1}j_{2}}=S_{N-1}\left( \hat{p}\right) \zeta _{p}^{\ast }\left\{
V_{f\,j_{1}p}^{\ast }V_{f\,j_{2}p}+V_{f\,j_{1}p+s}^{\ast
}V_{f\,j_{2}p+s}\right\}  \notag \\
&&\qquad \qquad \qquad \qquad \left. +\right. \sum_{j\neq
p}^{s}S_{N-2}\left( \hat{p}\hat{\jmath}\right) \zeta _{p}^{\ast
}n_{j}\left\{ V_{f\,j_{1}j}^{\ast }V_{f\,j_{2}j}+V_{f\,j_{1}j+s}^{\ast
}V_{f\,j_{2}j+s}\right\}
\end{eqnarray}%
and $TP_{2}$ is a term whose elements' computational cost scale as $s^{2}$
and is given by 
\begin{eqnarray}
TP_{2pj_{1}j_{2}j_{3}j_{4}} &=&S_{N-1}\left( \hat{p}\right) \zeta _{p}^{\ast
}V_{f\,j_{1}p}^{\ast }V_{f\,j_{2}p+s}^{\ast }V_{f\,j_{3}p}V_{f\,j_{4}p+s} 
\notag \\
&&+\sum_{j\neq p}^{s}S_{N-2}\left( \hat{p}\hat{\jmath}\right) \left\{ \zeta
_{p}^{\ast }n_{j}\left\{ V_{f\,j_{1}j}^{\ast }V_{f\,j_{2}j+s}^{\ast
}V_{f\,j_{3}j}V_{f\,j_{4}j+s}+V_{f\,j_{1}j}^{\ast }V_{f\,j_{2}p}^{\ast
}V_{f\,j_{3}p}V_{f\,j_{4}j}\right. \right.  \notag \\
&&\qquad +V_{f\,j_{1}j}^{\ast }V_{f\,j_{2}p+s}^{\ast
}V_{f\,j_{3}p+s}V_{f\,j_{4}j}+V_{f\,j_{1}j+s}^{\ast }V_{f\,j_{2}p}^{\ast
}V_{f\,j_{3}p}V_{f\,j_{4}j+s}  \notag \\
&&\qquad \qquad +\left. +V_{f\,j_{1}j+s}^{\ast }V_{f\,j_{2}p+s}^{\ast
}V_{f\,j_{3}p+s}V_{f\,j_{4}j+s}\right\} \left. +\,\zeta _{j}^{\ast
}V_{f\,j_{1}j}^{\ast }V_{f\,j_{2}j+s}^{\ast
}V_{f\,j_{3}p}V_{f\,j_{4}p+s}\right\}  \notag \\
&&+\sum_{j\neq k\neq p}^{s}S_{N-2}\left( \hat{p}\hat{\jmath}\hat{k}\right)
\zeta _{p}^{\ast }\zeta _{j}^{\ast }\zeta _{k}V_{f\,j_{1}j}^{\ast
}V_{f\,j_{2}j+s}^{\ast }V_{f\,j_{3}k}V_{f\,j_{4}k+s}  \notag \\
&&+\frac{1}{2}\zeta _{p}^{\ast }\sum_{j,k\neq p}^{s}S_{N-3}\left( \hat{p}%
\hat{\jmath}\hat{k}\right) n_{j}n_{k}\left\{ V_{f\,j_{1}j}^{\ast
}V_{f\,j_{2}k}^{\ast }V_{f\,j_{3}k}V_{f\,j_{4}pj}\right.  \notag \\
&&\qquad \qquad \qquad +V_{f\,j_{1}j}^{\ast }V_{f\,j_{2}k+s}^{\ast
}V_{f\,j_{3}k+s}V_{f\,j_{4}pj}+V_{f\,j_{1}j+s}^{\ast }V_{f\,j_{2}k}^{\ast
}V_{f\,j_{3}k}V_{f\,j_{4}pj+s}  \notag \\
&&\qquad \qquad \qquad \qquad \qquad \qquad \qquad \qquad \left.
+V_{f\,j_{1}j+s}^{\ast }V_{f\,j_{2}k+s}^{\ast
}V_{f\,j_{3}k+s}V_{f\,j_{4}pj+s}\right\}
\end{eqnarray}%
The unpaired term 
\begin{eqnarray}
&&\left\langle G^{\dagger N}\phi |Hb_{p}^{\dagger }b_{q}^{\dagger
}G^{\dagger N-1}\phi \right\rangle =N!\left( N-1\right) !\times  \notag \\
&&\left\{ S_{N-1}\left( \hat{p}\hat{q}\right)
\sum_{j_{1},j_{2}=1}^{r}\left\{ V_{f\,j_{1}\bar{q}}^{\ast
}V_{f\,j_{2}p}\zeta _{q}^{\ast }+V_{f\,j_{1}\bar{p}}^{\ast
}V_{f\,j_{2}q}\zeta _{p}^{\ast }\right\} h_{j_{1}j_{2}}^{a}\right.  \notag \\
&&+\sum_{j_{1,}j_{2,}j_{3,}j_{4}=1}^{r}\left\{ \sum_{\kappa \neq \left|
p\right| \neq \left| q\right| }^{s}S_{N-1}\left( \hat{\kappa}\hat{p}\hat{q}%
\right) \left\{ V_{f\,j_{1}\kappa }^{\ast }V_{f\,j_{2}\kappa +s}^{\ast
}V_{f\,j_{3}p}V_{f\,j_{4}q}\zeta _{\kappa }^{\ast }\right. \right.  \notag \\
&&\quad \qquad \quad +V_{f\,j_{1}\bar{p}}^{\ast }V_{f\,j_{2}\bar{q}}^{\ast
}V_{f\,j_{3}\kappa }V_{f\,j_{4}\kappa +s}\zeta _{p}^{\ast }\zeta _{q}^{\ast
}\zeta _{\kappa }+V_{f\,j_{1}\bar{q}}^{\ast }V_{f\,j_{2}\kappa }^{\ast
}V_{f\,j_{3}p}V_{f\,j_{4}\kappa }n_{\kappa }\zeta _{q}^{\ast }  \notag \\
&&\qquad \qquad \qquad \qquad \qquad \qquad \qquad \qquad \qquad \qquad
\left. +V_{f\,j_{1}\bar{p}}^{\ast }V_{f\,j_{2}\kappa }^{\ast
}V_{f\,j_{3}q}V_{f\,j_{4}\kappa }n_{\kappa }\zeta _{p}^{\ast }\right\} 
\notag \\
&&\left. +S_{N-1}\left( \hat{p}\hat{q}\right) \left\{ \zeta _{p}^{\ast
}V_{f\,j_{1}p}^{\ast }V_{f\,j_{2}p+s}^{\ast
}V_{f\,j_{3}p}V_{f\,j_{4}q}+\zeta _{q}^{\ast }V_{f\,j_{1}q}^{\ast
}V_{f\,j_{2}q+s}^{\ast }V_{f\,j_{3}p}V_{f\,j_{4}q}\right\} \right\}  \notag
\\
&&\qquad \qquad \qquad \qquad \qquad \qquad \qquad \qquad \qquad \qquad
\qquad \qquad \left. \left\langle j_{1}j_{2}||j_{3}j_{4}\right\rangle ^{a}%
\overset{}{\underset{}{}}\right\}
\end{eqnarray}%
can be expressed as a sum of two types $TU_{1}$ and $TU_{2}$ 
\begin{eqnarray}
&&\left\langle G^{\dagger N}\phi |Hb_{p}^{\dagger }b_{q}^{\dagger
}G^{\dagger N-1}\phi \right\rangle =TU_{1}+TU_{2}  \notag \\
&&\left. {}\right. =N!\left( N-1\right) !\left\{
\sum_{j_{1},j_{2}=1}^{r}TU_{1pqj_{1}j_{2}}h_{j_{1}j_{2}}^{a}+%
\sum_{j_{1,}j_{2,}j_{3,}j_{4}=1}^{r}TU_{2pqj_{1}j_{2}j_{3}j_{4}}\left\langle
j_{1}j_{2}||j_{3}j_{4}\right\rangle ^{a}\right\}
\end{eqnarray}%
where $TU_{1}$ is given by 
\begin{equation}
TU_{1pqj_{1}j_{2}}=S_{N-1}\left( \hat{p}\hat{q}\right) \left\{ V_{f\,j_{1}%
\bar{q}}^{\ast }V_{f\,j_{2}p}\zeta _{q}^{\ast }+V_{f\,j_{1}\bar{p}}^{\ast
}V_{f\,j_{2}q}\zeta _{p}^{\ast }\right\}
\end{equation}%
and $TU_{2}$ is a term whose elements' computational cost scale as $s$ and
is given by 
\begin{eqnarray}
TU_{2pqj_{1}j_{2}j_{3}j_{4}} &=&\sum_{\kappa \neq \left| p\right| \neq
\left| q\right| }^{s}S_{N-1}\left( \hat{\kappa}\hat{p}\hat{q}\right) \left\{
V_{f\,j_{1}\kappa }^{\ast }V_{f\,j_{2}\kappa +s}^{\ast
}V_{f\,j_{3}p}V_{f\,j_{4}q}\zeta _{\kappa }^{\ast }\right.  \notag \\
&&\qquad +V_{f\,j_{1}\bar{p}}^{\ast }V_{f\,j_{2}\bar{q}}^{\ast
}V_{f\,j_{3}\kappa }V_{f\,j_{4}\kappa +s}\zeta _{p}^{\ast }\zeta _{q}^{\ast
}\zeta _{\kappa }+V_{f\,j_{1}\bar{q}}^{\ast }V_{f\,j_{2}\kappa }^{\ast
}V_{f\,j_{3}p}V_{f\,j_{4}\kappa }n_{\kappa }\zeta _{q}^{\ast }  \notag \\
&&\qquad \qquad \qquad \qquad \qquad \qquad \qquad \qquad \left. +V_{f\,j_{1}%
\bar{p}}^{\ast }V_{f\,j_{2}\kappa }^{\ast }V_{f\,j_{3}q}V_{f\,j_{4}\kappa
}n_{\kappa }\zeta _{p}^{\ast }\right\}  \notag \\
&&+S_{N-1}\left( \hat{p}\hat{q}\right) \left\{ \zeta _{p}^{\ast
}V_{f\,j_{1}p}^{\ast }V_{f\,j_{2}p+s}^{\ast
}V_{f\,j_{3}p}V_{f\,j_{4}q}+\zeta _{q}^{\ast }V_{f\,j_{1}q}^{\ast
}V_{f\,j_{2}q+s}^{\ast }V_{f\,j_{3}p}V_{f\,j_{4}q}\right\}
\end{eqnarray}

The partial derivative $\left\langle G^{\dagger N}\phi |Ha_{i_{1}}^{\dagger
}a_{i_{2}}^{\dagger }G^{\dagger N-1}\phi \right\rangle $ is the sum of
paired and unpaired terms given by 
\begin{eqnarray}
&&\left\langle G^{\dagger N}\phi |Ha_{i_{1}}^{\dagger }a_{i_{2}}^{\dagger
}G^{\dagger N-1}\phi \right\rangle =\sum_{p=1}^{r}\left\{
V_{f\,i_{1}p}V_{f\,i_{2}p+s}-V_{f\,i_{1}p+s}V_{f\,i_{2}p}\right\}
\left\langle G^{\dagger N}\phi |Hb_{p}^{\dagger }b_{p+s}^{\dagger
}G^{\dagger N-1}\phi \right\rangle  \notag \\
&&\qquad \qquad \qquad \qquad \qquad \qquad +\sum_{\substack{ p<q=1  \\ %
q\neq p+s}}^{r}\left\{
V_{f\,i_{1}p}V_{f\,j_{2}q}-V_{f\,i_{1}q}V_{f\,i_{2}p}\right\} \left\langle
G^{\dagger N}\phi |Hb_{p}^{\dagger }b_{q}^{\dagger }G^{\dagger N-1}\phi
\right\rangle
\end{eqnarray}%
which can be expressed in terms of $TP_{0}$, $TP_{1}$, $TP_{2}$, $TU_{1}$
and $TU_{2}$ as 
\begin{eqnarray}
&&\left\langle G^{\dagger N}\phi |Ha_{i_{1}}^{\dagger }a_{i_{2}}^{\dagger
}G^{\dagger N-1}\phi \right\rangle =N!\left( N-1\right) !\left\{
\sum_{p=1}^{r}\left\{
V_{f\,i_{1}p}V_{f\,i_{2}p+s}-V_{f\,i_{1}p+s}V_{f\,i_{2}p}\right\} \times
\right.  \notag \\
&&\qquad \qquad \qquad \qquad \left\{
\sum_{j_{1,}j_{2}=1}^{r}TP_{1pj_{1}j_{2}}h_{j_{1}j_{2}}^{a}+%
\sum_{j_{1,}j_{2,}j_{3,}j_{4}=1}^{r}TP_{2pj_{1}j_{2}j_{3}j_{4}}\left\langle
j_{1}j_{2}||j_{3}j_{4}\right\rangle ^{a}\right\} +  \notag \\
&&\sum_{\substack{ p<q=1  \\ q\neq p+s}}^{r}\left\{
V_{f\,i_{1}p}V_{f\,j_{2}q}-V_{f\,i_{1}q}V_{f\,i_{2}p}\right\} \left\{
\sum_{j_{1},j_{2}=1}^{r}TU_{1pqj_{1}j_{2}}h_{j_{1}j_{2}}^{a}+%
\sum_{j_{1,}j_{2,}j_{3,}j_{4}=1}^{r}TU_{2pqj_{1}j_{2}j_{3}j_{4}}\left\langle
j_{1}j_{2}||j_{3}j_{4}\right\rangle ^{a}\right\}  \notag \\
&&
\end{eqnarray}%
that can be written as 
\begin{equation}
\left\langle G^{\dagger N}\phi |Ha_{i_{1}}^{\dagger }a_{i_{2}}^{\dagger
}G^{\dagger N-1}\phi \right\rangle =\mathcal{D}_{0i_{1}i_{2}}h_{0}+%
\sum_{j_{1,}j_{2}=1}^{r}\mathcal{D}%
_{1i_{1}i_{2}j_{1}j_{2}}h_{j_{1}j_{2}}^{a}+%
\sum_{j_{1,}j_{2,}j_{3,}j_{4}=1}^{r}\mathcal{D}%
_{2i_{1}i_{2}j_{1}j_{2}j_{3}j_{4}}\left\langle
j_{1}j_{2}||j_{3}j_{4}\right\rangle ^{a}  \label{matel}
\end{equation}%
where $\mathcal{D}_{0i_{1}i_{2}}$ is a term whose elements' computational
cost scale as $s$ and is given by 
\begin{equation}
\mathcal{D}_{0i_{1}i_{2}}=N!\left( N-1\right) !\sum_{p=1}^{s}\left\{
V_{f\,i_{1}p}V_{f\,i_{2}p+s}-V_{f\,i_{1}p+s}V_{f\,i_{2}p}\right\} TP_{0p}
\end{equation}%
$\mathcal{D}_{1i_{1}i_{2}}$ is a term whose elements' computation scales as $%
s^{3}$ given by 
\begin{eqnarray}
\mathcal{D}_{1i_{1}i_{2}j_{1}j_{2}} &=&N!\left( N-1\right) !\left\{
\sum_{p=1}^{s}\left\{
V_{f\,i_{1}p}V_{f\,i_{2}p+s}-V_{f\,i_{1}p+s}V_{f\,i_{2}p}\right\}
TP_{1pj_{1}j_{2}}\right.  \notag \\
&&\qquad \qquad \qquad \qquad \left. +\sum_{\substack{ p<q=1  \\ q\neq p+s}}%
^{s}\left\{ V_{f\,i_{1}p}V_{f\,j_{2}q}-V_{f\,i_{1}q}V_{f\,i_{2}p}\right\}
TU_{1pqj_{1}j_{2}}\right\}
\end{eqnarray}%
and $\mathcal{D}_{2i_{1}i_{2}}$ is a term whose elements' computational cost
scale as $s^{4}$ and is given by%
\begin{eqnarray}
\mathcal{D}_{2i_{1}i_{2}j_{1}j_{2}j_{3}j_{4}} &=&N!\left( N-1\right)
!\left\{ \sum_{p=1}^{s}\left\{
V_{f\,i_{1}p}V_{f\,i_{2}p+s}-V_{f\,i_{1}p+s}V_{f\,i_{2}p}\right\}
TP_{2pj_{1}j_{2}j_{3}j_{4}}\right.  \notag \\
&&\qquad \qquad \qquad +\left. \sum_{\substack{ p<q=1  \\ q\neq p+s}}%
^{s}\left\{ V_{f\,i_{1}p}V_{f\,j_{2}q}-V_{f\,i_{1}q}V_{f\,i_{2}p}\right\}
TU_{2pqj_{1}j_{2}j_{3}j_{4}}\right\}
\end{eqnarray}

The calculation of $TP_{0}$, $TP_{1}$, $TP_{2}$, $TU_{1}$ and $TU_{2}$ need
the following elementary symmetric functions 
\begin{equation}
S_{N-1}\left( \hat{p}\right) ,S_{N-1}\left( \hat{p}\hat{q}\right)
,S_{N-1}\left( \hat{\kappa}\hat{p}\hat{q}\right) ,S_{N-2}\left( \hat{p}\hat{%
\jmath}\right) ,S_{N-2}\left( \hat{p}\hat{\jmath}\hat{k}\right)
,S_{N-3}\left( \hat{p}\hat{\jmath}\hat{k}\right)
\end{equation}%
to be evaluated. These can be computed independently of the other terms in
the calculation by a very efficient, robust procedure that scales as $s^{4}$%
, which has been developed by \cite{Jensen80,Jensen81}. This procedure is
based on a graphical approach to a recursion relationship between the
symmetric functions $S_{N},S_{N-1}\left( \hat{p}\right) $ and $S_{N}\left( 
\hat{p}\right) $ described in these references.

The number of steps in the computation of $\left\langle G^{\dagger N}\phi
|Ha_{i_{1}}^{\dagger }a_{i_{2}}^{\dagger }G^{\dagger N-1}\phi \right\rangle $
in eq. $\left( \ref{matel}\right) $ is $r^{5}$, given that $\left\{ \mathcal{%
D}_{2i_{1}i_{2}j_{1}j_{2}j_{3}j_{4}}\right\} $ and $\left\{ \left\langle
j_{1}j_{2}||j_{3}j_{4}\right\rangle ^{a}\right\} $ have been calculated,
thus form the rate controlling step and the overall computational cost of
the derivatives scale as $r^{5}$.

\bibliographystyle{apsrev}
\bibliography{AGPOPT,BW}

\end{document}
